% Options for packages loaded elsewhere
\PassOptionsToPackage{unicode}{hyperref}
\PassOptionsToPackage{hyphens}{url}
\documentclass[
  nobib]{tufte-handout}
\usepackage{xcolor}
\usepackage{amsmath,amssymb}
\setcounter{secnumdepth}{-\maxdimen} % remove section numbering
\usepackage{iftex}
\ifPDFTeX
  \usepackage[T1]{fontenc}
  \usepackage[utf8]{inputenc}
  \usepackage{textcomp} % provide euro and other symbols
\else % if luatex or xetex
  \usepackage{unicode-math} % this also loads fontspec
  \defaultfontfeatures{Scale=MatchLowercase}
  \defaultfontfeatures[\rmfamily]{Ligatures=TeX,Scale=1}
\fi
\usepackage{lmodern}
\ifPDFTeX\else
  % xetex/luatex font selection
\fi
% Use upquote if available, for straight quotes in verbatim environments
\IfFileExists{upquote.sty}{\usepackage{upquote}}{}
\IfFileExists{microtype.sty}{% use microtype if available
  \usepackage[]{microtype}
  \UseMicrotypeSet[protrusion]{basicmath} % disable protrusion for tt fonts
}{}
\makeatletter
\@ifundefined{KOMAClassName}{% if non-KOMA class
  \IfFileExists{parskip.sty}{%
    \usepackage{parskip}
  }{% else
    \setlength{\parindent}{0pt}
    \setlength{\parskip}{6pt plus 2pt minus 1pt}}
}{% if KOMA class
  \KOMAoptions{parskip=half}}
\makeatother
\usepackage{color}
\usepackage{fancyvrb}
\newcommand{\VerbBar}{|}
\newcommand{\VERB}{\Verb[commandchars=\\\{\}]}
\DefineVerbatimEnvironment{Highlighting}{Verbatim}{commandchars=\\\{\}}
% Add ',fontsize=\small' for more characters per line
\newenvironment{Shaded}{}{}
\newcommand{\AlertTok}[1]{\textcolor[rgb]{1.00,0.00,0.00}{\textbf{#1}}}
\newcommand{\AnnotationTok}[1]{\textcolor[rgb]{0.38,0.63,0.69}{\textbf{\textit{#1}}}}
\newcommand{\AttributeTok}[1]{\textcolor[rgb]{0.49,0.56,0.16}{#1}}
\newcommand{\BaseNTok}[1]{\textcolor[rgb]{0.25,0.63,0.44}{#1}}
\newcommand{\BuiltInTok}[1]{\textcolor[rgb]{0.00,0.50,0.00}{#1}}
\newcommand{\CharTok}[1]{\textcolor[rgb]{0.25,0.44,0.63}{#1}}
\newcommand{\CommentTok}[1]{\textcolor[rgb]{0.38,0.63,0.69}{\textit{#1}}}
\newcommand{\CommentVarTok}[1]{\textcolor[rgb]{0.38,0.63,0.69}{\textbf{\textit{#1}}}}
\newcommand{\ConstantTok}[1]{\textcolor[rgb]{0.53,0.00,0.00}{#1}}
\newcommand{\ControlFlowTok}[1]{\textcolor[rgb]{0.00,0.44,0.13}{\textbf{#1}}}
\newcommand{\DataTypeTok}[1]{\textcolor[rgb]{0.56,0.13,0.00}{#1}}
\newcommand{\DecValTok}[1]{\textcolor[rgb]{0.25,0.63,0.44}{#1}}
\newcommand{\DocumentationTok}[1]{\textcolor[rgb]{0.73,0.13,0.13}{\textit{#1}}}
\newcommand{\ErrorTok}[1]{\textcolor[rgb]{1.00,0.00,0.00}{\textbf{#1}}}
\newcommand{\ExtensionTok}[1]{#1}
\newcommand{\FloatTok}[1]{\textcolor[rgb]{0.25,0.63,0.44}{#1}}
\newcommand{\FunctionTok}[1]{\textcolor[rgb]{0.02,0.16,0.49}{#1}}
\newcommand{\ImportTok}[1]{\textcolor[rgb]{0.00,0.50,0.00}{\textbf{#1}}}
\newcommand{\InformationTok}[1]{\textcolor[rgb]{0.38,0.63,0.69}{\textbf{\textit{#1}}}}
\newcommand{\KeywordTok}[1]{\textcolor[rgb]{0.00,0.44,0.13}{\textbf{#1}}}
\newcommand{\NormalTok}[1]{#1}
\newcommand{\OperatorTok}[1]{\textcolor[rgb]{0.40,0.40,0.40}{#1}}
\newcommand{\OtherTok}[1]{\textcolor[rgb]{0.00,0.44,0.13}{#1}}
\newcommand{\PreprocessorTok}[1]{\textcolor[rgb]{0.74,0.48,0.00}{#1}}
\newcommand{\RegionMarkerTok}[1]{#1}
\newcommand{\SpecialCharTok}[1]{\textcolor[rgb]{0.25,0.44,0.63}{#1}}
\newcommand{\SpecialStringTok}[1]{\textcolor[rgb]{0.73,0.40,0.53}{#1}}
\newcommand{\StringTok}[1]{\textcolor[rgb]{0.25,0.44,0.63}{#1}}
\newcommand{\VariableTok}[1]{\textcolor[rgb]{0.10,0.09,0.49}{#1}}
\newcommand{\VerbatimStringTok}[1]{\textcolor[rgb]{0.25,0.44,0.63}{#1}}
\newcommand{\WarningTok}[1]{\textcolor[rgb]{0.38,0.63,0.69}{\textbf{\textit{#1}}}}
\usepackage{longtable,booktabs,array}
\newcounter{none} % for unnumbered tables
\usepackage{calc} % for calculating minipage widths
% Correct order of tables after \paragraph or \subparagraph
\usepackage{etoolbox}
\makeatletter
\patchcmd\longtable{\par}{\if@noskipsec\mbox{}\fi\par}{}{}
\makeatother
% Allow footnotes in longtable head/foot
\IfFileExists{footnotehyper.sty}{\usepackage{footnotehyper}}{\usepackage{footnote}}
\makesavenoteenv{longtable}
\setlength{\emergencystretch}{3em} % prevent overfull lines
\providecommand{\tightlist}{%
  \setlength{\itemsep}{0pt}\setlength{\parskip}{0pt}}
% =====================================================================
% axona-doc-preamble.tex — shared LaTeX preamble for the programmer-guide
% quartet (Quick-Start, Programmer Guide, API Reference, Services Guide).
%
% These docs are markdown-native; the PDF is built by converting the .md to
% a tufte-handout .tex with pandoc (-H this file) and compiling with tectonic.
% See render.sh. Modeled on explainer/Axona-Explainer.tex (the flat-essay
% original); the tweaks below make tufte-handout work for code/table-heavy
% reference material.
% =====================================================================

\definecolor{rust}{RGB}{185,78,54}
\definecolor{rustsoft}{RGB}{212,168,150}
\definecolor{inkmuted}{RGB}{102,102,102}
\usepackage{microtype}

% Explainer-matched page furniture: tufte-handout's letterspaced running
% heads fail on certain page breaks (the SOUL \MakeTextLowercase issue the
% explainer works around), so use the same fancyhdr style — no running
% head, a centered page number in the footer.
\usepackage{fancyhdr}
\fancypagestyle{ndhtplain}{%
  \fancyhf{}%
  \renewcommand{\headrulewidth}{0pt}%
  \fancyfoot[C]{\thepage}%
}

% Explainer-matched heading-orphan protection: reserve vertical room before
% each (sub)section heading so a heading never sits alone at a page bottom.
\usepackage{etoolbox}
\usepackage{needspace}
\pretocmd{\section}{\needspace{6\baselineskip}}{}{}
\pretocmd{\subsection}{\needspace{4\baselineskip}}{}{}

% XeTeX (tectonic) unicode monospace: renders … → and the box-drawing ASCII
% diagrams that the default CM/LM mono cannot. Menlo ships with macOS.
\setmonofont{Menlo}[Scale=0.78]

% tufte-handout ships \subsubsection as an error stub (it wants flat essays).
% A run-in fallback lets nested reference headings (§4.2, §12.3, …) compile.
\renewcommand{\subsubsection}[1]{\medskip\par\noindent{\normalsize\bfseries #1}\par\nobreak\smallskip}

% These docs use no margin notes, so reclaim tufte's wide right margin: widen
% the text block so code blocks and API tables don't overflow it.
\usepackage{geometry}
\geometry{left=1in, textwidth=6.1in, marginparsep=0.2in, marginparwidth=0.9in,
          top=1in, headsep=0.4in}

% Wrap long code lines (cheat-sheet comments run to ~130 chars) with a ↪ marker
% instead of overflowing. Applies to pandoc's fancyvrb Highlighting/Verbatim.
\usepackage{fvextra}
\fvset{breaklines=true, breakanywhere=true}

% A few technical symbols appear in prose / inline code that Latin Modern and
% Menlo don't cover; map them to math equivalents so they render (rather than
% dropping to tofu). Only maps chars actually used — leaves → alone (it renders).
\usepackage{newunicodechar}
\newunicodechar{≈}{\ensuremath{\approx}}
\newunicodechar{≥}{\ensuremath{\geq}}
\newunicodechar{≤}{\ensuremath{\leq}}
\newunicodechar{×}{\ensuremath{\times}}
\newunicodechar{·}{\ensuremath{\cdot}}
\newunicodechar{‖}{\ensuremath{\|}}
\usepackage{bookmark}
\IfFileExists{xurl.sty}{\usepackage{xurl}}{} % add URL line breaks if available
\urlstyle{same}
\hypersetup{
  hidelinks,
  pdfcreator={LaTeX via pandoc}}

\author{}
\date{}

\begin{document}

% Generated by render.sh — the explainer-style title page.
\thispagestyle{empty}
\null\vfill
\begin{center}
{\fontsize{48}{52}\selectfont \textsc{Axona}\par}
\vspace{0.6em}
{\Large\itshape Every Public Symbol, Exactly\par}
\vspace{4em}
{\large David A.~Smith\par}
{\itshape\small Axona.net\par}
\vspace{2em}
{\small The Axona API Reference \quad v4.22.0 \quad 2026-07-14\par}
\end{center}
\vfill
\null
\newpage
\pagestyle{ndhtplain}
\thispagestyle{ndhtplain}
\setcounter{page}{1}

\section{Axona API Reference}\label{axona-api-reference}

Reference for every public symbol exported from \texttt{@axona/protocol}
v4.22.0.

Organized by what application developers reach for first (identity, peer
lifecycle, pub/sub, direct messaging, introspection), then the
transport/protocol surface, then low-level utilities. Every signature
below is verified against the v4.22.0 kernel source.

\begin{quote}
\textbf{Which network?} Both live networks run the 4.x line (wire 4.0).
\textbf{Testnet} (\texttt{wss://testnet.axona.net}) tracks the newest
kernel --- v4.22.0, the version this reference describes.
\textbf{Production} (\texttt{wss://bridge.axona.net} east +
\texttt{wss://bridge-west.axona.net} west) runs the most recently
promoted kernel, typically one release behind while changes soak. The
API surface below is identical on both; point wherever you deploy.
Install \texttt{github:axona-net/axona-protocol\#v4.22.0}.
\end{quote}

Companion documents:

\begin{itemize}
\tightlist
\item
  \href{Quick-Start-v4.22.0.md}{Quick Start} --- 5-minute working
  roundtrip.
\item
  \href{Axona-Programmer-Guide-v4.22.0.md}{Programmer Guide} --- mental
  model + worked example + pitfalls.
\item
  \href{../SECURITY-CHANGELOG.md}{Security changelog} --- what each
  kernel version protects.
\end{itemize}

Imports throughout assume:

\begin{Shaded}
\begin{Highlighting}[]
\ImportTok{import}\NormalTok{ \{ }\CommentTok{/* … */}\NormalTok{ \} }\ImportTok{from} \StringTok{\textquotesingle{}@axona/protocol\textquotesingle{}}\OperatorTok{;}
\end{Highlighting}
\end{Shaded}

Sub-path imports (e.g.~\texttt{@axona/protocol/contracts/Transport.js},
\texttt{@axona/protocol/persistence/file.js}) work for symbols that need
them, but most apps only need the main barrel.

\begin{quote}
\textbf{How to read this reference.} It has three parts. The
\textbf{Application surface} (§§1--9) is what applications call --- most
readers never need anything else. The \textbf{Transport + protocol
surface} (§§10--16) and \textbf{Low-level utilities} (§§17--23) exist
for transport authors, tooling, and the curious. Building with an AI
assistant? Hand it the \href{Axona-AI-Grounding-v4.22.0.md}{AI
Grounding} file --- the application surface distilled to rules and
patterns.
\end{quote}

\begin{center}\rule{0.5\linewidth}{0.5pt}\end{center}

\subsection{Table of contents}\label{table-of-contents}

\subsubsection{Application surface (what you'll use
first)}\label{application-surface-what-youll-use-first}

\begin{enumerate}
\def\labelenumi{\arabic{enumi}.}
\tightlist
\item
  \hyperref[1-types-referenced-throughout]{Types referenced throughout}
\item
  \hyperref[2-identity]{Identity}
\item
  \hyperref[3-axonapeer-construction--lifecycle]{AxonaPeer construction
  + lifecycle}
\item
  \hyperref[4-pubsub]{Pub/sub}

  \begin{itemize}
  \tightlist
  \item
    \hyperref[41-topic-addressing]{4.1 Topic addressing}
  \item
    \hyperref[42-publish--subscribe]{4.2 publish / subscribe}
  \item
    \hyperref[43-pull--metrics]{4.3 pull / metrics}
  \item
    \hyperref[44-owner--creator-ops-kill]{4.4 owner + creator ops: kill}
  \item
    \hyperref[45-hosting-host--unhost]{4.5 hosting: host / unhost}
  \item
    \hyperref[46-topic-limits]{4.6 topic limits}
  \end{itemize}
\item
  \hyperref[5-direct-messaging]{Direct messaging}
\item
  \hyperref[6-mesh-introspection--events]{Mesh introspection + events}
\item
  \hyperref[7-envelopes--verification]{Envelopes + verification}
\item
  \hyperref[8-errors]{Errors}
\item
  \hyperref[9-persistence]{Persistence}
\end{enumerate}

\subsubsection{Transport + protocol
surface}\label{transport-protocol-surface}

\begin{quote}
\textbf{You probably don't need this part.} Everything an application
calls is in §§1--9 above. The sections from here on are for people
implementing transports, embedding the kernel, or building tooling.
\end{quote}

\begin{enumerate}
\def\labelenumi{\arabic{enumi}.}
\setcounter{enumi}{9}
\tightlist
\item
  \hyperref[10-contracts]{Contracts}
\item
  \hyperref[11-sim-transport]{Sim transport}
\item
  \hyperref[12-web-transport]{Web transport}
\item
  \hyperref[13-handshake]{Handshake}
\item
  \hyperref[14-authenticated-identity-handshake-axona5]{Authenticated-identity
  handshake (axona/5)}
\item
  \hyperref[15-bridge-directory]{Bridge directory}
\item
  \hyperref[16-low-level-pubsub-builders]{Low-level pub/sub builders}
\end{enumerate}

\subsubsection{Low-level utilities}\label{low-level-utilities}

\begin{quote}
\textbf{Internals.} Pure helpers the kernel itself is built from. Apps
import these only for unusual jobs (custom signing, ID math, region
tables).
\end{quote}

\begin{enumerate}
\def\labelenumi{\arabic{enumi}.}
\setcounter{enumi}{16}
\tightlist
\item
  \hyperref[17-ed25519-helpers]{Ed25519 helpers}
\item
  \hyperref[18-hex--id-math]{Hex / ID math}
\item
  \hyperref[19-s2-geographic-cells]{S2 geographic cells}
\item
  \hyperref[20-region-names]{Region names}
\item
  \hyperref[21-geo-helpers]{Geo helpers}
\item
  \hyperref[22-proof-of-work-scaffolding]{Proof-of-work scaffolding}
\item
  \hyperref[23-constants]{Constants}
\end{enumerate}

\begin{center}\rule{0.5\linewidth}{0.5pt}\end{center}

\subsection{Application surface}\label{application-surface}

\subsection{1. Types referenced
throughout}\label{types-referenced-throughout}

\subsubsection{\texorpdfstring{\texttt{Identity} (node
identity)}{Identity (node identity)}}\label{identity-type}

Returned by \texttt{createNodeIdentity}. The connection/transport
keypair; its pubkey forms the nodeId.

\begin{Shaded}
\begin{Highlighting}[]
\NormalTok{\{}
\NormalTok{  id}\OperatorTok{:}         \DataTypeTok{string}\OperatorTok{;}       \CommentTok{// 66{-}char lowercase hex (264{-}bit nodeId)}
\NormalTok{  pubkey}\OperatorTok{:}     \BuiltInTok{Uint8Array}\OperatorTok{;}   \CommentTok{// 32{-}byte Ed25519 public key}
\NormalTok{  pubkeyHex}\OperatorTok{:}  \DataTypeTok{string}\OperatorTok{;}       \CommentTok{// 64{-}char lowercase hex}
\NormalTok{  privateKey}\OperatorTok{:} \BuiltInTok{CryptoKey}\OperatorTok{;}    \CommentTok{// Web Crypto Ed25519 (browser + Node 20+)}
\NormalTok{  region}\OperatorTok{:}\NormalTok{     \{ lat}\OperatorTok{:} \DataTypeTok{number}\OperatorTok{,}\NormalTok{ lng}\OperatorTok{:} \DataTypeTok{number}\NormalTok{ \}}\OperatorTok{;}
\NormalTok{  createdAt}\OperatorTok{:}  \DataTypeTok{number}\OperatorTok{;}       \CommentTok{// ms epoch}
\NormalTok{  pow}\OperatorTok{:}        \DataTypeTok{string}\OperatorTok{;}       \CommentTok{// transport{-}role PoW nonce (\textquotesingle{}\textquotesingle{} = inert at difficulty 0)}
  \FunctionTok{sign}\NormalTok{(message}\OperatorTok{:} \BuiltInTok{Uint8Array}\NormalTok{)}\OperatorTok{:} \BuiltInTok{Promise}\OperatorTok{\textless{}}\BuiltInTok{Uint8Array}\OperatorTok{\textgreater{};}
  \FunctionTok{verify}\NormalTok{(message}\OperatorTok{:} \BuiltInTok{Uint8Array}\OperatorTok{,}\NormalTok{ signature}\OperatorTok{:} \BuiltInTok{Uint8Array}\NormalTok{)}\OperatorTok{:} \BuiltInTok{Promise}\OperatorTok{\textless{}}\DataTypeTok{boolean}\OperatorTok{\textgreater{};}
\NormalTok{\}}
\end{Highlighting}
\end{Shaded}

The top byte of \texttt{id} is the S2 region cell for
\texttt{(lat,\ lng)}; the bottom 256 bits are \texttt{SHA-256(pubkey)}.
The id is self-authenticating --- a peer can only claim an id it holds
the private key for.

\subsubsection{\texorpdfstring{\texttt{AuthorIdentity} (authorship
identity)}{AuthorIdentity (authorship identity)}}\label{author-identity-type}

Returned by \texttt{createAuthorIdentity}. A keypair only --- \textbf{no
nodeId, no region} (authorship is not a place). Its public key is the
Author ID.

\begin{Shaded}
\begin{Highlighting}[]
\NormalTok{\{}
\NormalTok{  kind}\OperatorTok{:}       \StringTok{\textquotesingle{}author\textquotesingle{}}\OperatorTok{;}
\NormalTok{  authorId}\OperatorTok{:}   \DataTypeTok{string}\OperatorTok{;}       \CommentTok{// 64{-}char hex Author ID (== signerPubkey on the wire)}
\NormalTok{  pubkey}\OperatorTok{:}     \BuiltInTok{Uint8Array}\OperatorTok{;}   \CommentTok{// 32{-}byte Ed25519 public key}
\NormalTok{  pubkeyHex}\OperatorTok{:}  \DataTypeTok{string}\OperatorTok{;}       \CommentTok{// 64{-}char hex (same value as authorId)}
\NormalTok{  privateKey}\OperatorTok{:} \BuiltInTok{CryptoKey}\OperatorTok{;}
\NormalTok{  createdAt}\OperatorTok{:}  \DataTypeTok{number}\OperatorTok{;}
\NormalTok{  pow}\OperatorTok{:}        \DataTypeTok{string}\OperatorTok{;}       \CommentTok{// publish{-}role PoW nonce (\textquotesingle{}\textquotesingle{} = inert)}
  \FunctionTok{sign}\NormalTok{(message}\OperatorTok{:} \BuiltInTok{Uint8Array}\NormalTok{)}\OperatorTok{:} \BuiltInTok{Promise}\OperatorTok{\textless{}}\BuiltInTok{Uint8Array}\OperatorTok{\textgreater{};}
  \FunctionTok{verify}\NormalTok{(message}\OperatorTok{:} \BuiltInTok{Uint8Array}\OperatorTok{,}\NormalTok{ signature}\OperatorTok{:} \BuiltInTok{Uint8Array}\NormalTok{)}\OperatorTok{:} \BuiltInTok{Promise}\OperatorTok{\textless{}}\DataTypeTok{boolean}\OperatorTok{\textgreater{};}
\NormalTok{\}}
\end{Highlighting}
\end{Shaded}

You pass an \texttt{AuthorIdentity} (or \texttt{ANONYMOUS}) as
\texttt{\{\ signWith\ \}} to \texttt{peer.pub} / \texttt{kill}.

\subsubsection{\texorpdfstring{\texttt{TopicDescriptor}
(type)}{TopicDescriptor (type)}}\label{topic-descriptor-type}

\begin{Shaded}
\begin{Highlighting}[]
\NormalTok{\{}
\NormalTok{  region}\OperatorTok{?:} \DataTypeTok{string} \OperatorTok{|} \DataTypeTok{number}\OperatorTok{;}   \CommentTok{// region name (\textquotesingle{}useast\textquotesingle{}), \textquotesingle{}0x89\textquotesingle{}, \textquotesingle{}137\textquotesingle{}, or 137; omit → node region}
\NormalTok{  owner}\OperatorTok{?:}  \DataTypeTok{string} \OperatorTok{|} \DataTypeTok{null}\OperatorTok{;}     \CommentTok{// 64{-}hex Author ID, or absent for an open topic}
\NormalTok{  name}\OperatorTok{:}    \DataTypeTok{string}\OperatorTok{;}            \CommentTok{// required, non{-}empty}
\NormalTok{  write}\OperatorTok{?:}  \StringTok{\textquotesingle{}open\textquotesingle{}} \OperatorTok{|} \StringTok{\textquotesingle{}owner\textquotesingle{}}\OperatorTok{;}  \CommentTok{// defaults by owner presence (see §4.1)}
\NormalTok{\}}
\end{Highlighting}
\end{Shaded}

\subsubsection{\texorpdfstring{\texttt{Envelope}
(type)}{Envelope (type)}}\label{envelope-type}

What \texttt{peer.sub} delivers and \texttt{peer.pull} returns.

\begin{Shaded}
\begin{Highlighting}[]
\NormalTok{\{}
\NormalTok{  msgId}\OperatorTok{:}        \DataTypeTok{string}\OperatorTok{;}   \CommentTok{// 64{-}char hex = sha256(canonical(\{ publisher, message \}))}
\NormalTok{  seq}\OperatorTok{:}          \DataTypeTok{number}\OperatorTok{;}   \CommentTok{// per{-}publisher monotonic sequence (folded under the signature)}
\NormalTok{  ts}\OperatorTok{:}           \DataTypeTok{number}\OperatorTok{;}   \CommentTok{// ms epoch when published}
\NormalTok{  topic}\OperatorTok{:}\NormalTok{        TopicDescriptor}\OperatorTok{;}   \CommentTok{// the SIGNED descriptor \{ region, owner, name, write \}}
\NormalTok{  message}\OperatorTok{:}      \DataTypeTok{any}\OperatorTok{;}      \CommentTok{// your JSON{-}serializable payload}
\NormalTok{  signature}\OperatorTok{?:}   \DataTypeTok{string}\OperatorTok{;}   \CommentTok{// \textquotesingle{}ed25519:\textless{}128 hex\textgreater{}\textquotesingle{}; present iff signed}
\NormalTok{  signerPubkey}\OperatorTok{?:} \DataTypeTok{string}\OperatorTok{;}  \CommentTok{// 64{-}char hex Author ID; present iff signed}
\NormalTok{  signerPow}\OperatorTok{?:}   \DataTypeTok{string}\OperatorTok{;}   \CommentTok{// publish{-}role PoW nonce; present iff signed (\textquotesingle{}\textquotesingle{} at difficulty 0)}
\NormalTok{\}}
\end{Highlighting}
\end{Shaded}

A \textbf{retraction} (see \texttt{peer.kill}) is delivered to your
\texttt{sub} handler as a marker instead of a full envelope:

\begin{Shaded}
\begin{Highlighting}[]
\NormalTok{\{ deleted}\OperatorTok{:} \KeywordTok{true}\OperatorTok{,}\NormalTok{ msgId}\OperatorTok{:} \DataTypeTok{string}\OperatorTok{,}\NormalTok{ topic}\OperatorTok{:} \DataTypeTok{string} \OperatorTok{|} \DataTypeTok{null}\NormalTok{ \}}
\end{Highlighting}
\end{Shaded}

Branch on \texttt{env.deleted\ ===\ true} to drop your local copy of
\texttt{msgId}.

\subsubsection{\texorpdfstring{\texttt{Subscription}
(type)}{Subscription (type)}}\label{subscription-type}

Opaque handle returned by \texttt{peer.sub()}.

\begin{Shaded}
\begin{Highlighting}[]
\NormalTok{sub}\OperatorTok{.}\AttributeTok{id}          \CommentTok{// string — unique per subscription}
\NormalTok{sub}\OperatorTok{.}\AttributeTok{topicId}     \CommentTok{// string — 66{-}char hex topic ID}
\NormalTok{sub}\OperatorTok{.}\AttributeTok{topicName}   \CommentTok{// string — your original topic name (or \textquotesingle{}\#\textless{}id{-}prefix\textgreater{}\textquotesingle{} for an id{-}only sub)}
\NormalTok{sub}\OperatorTok{.}\FunctionTok{stop}\NormalTok{()      }\CommentTok{// Promise\textless{}void\textgreater{} — cancel; idempotent}
\end{Highlighting}
\end{Shaded}

\subsubsection{\texorpdfstring{\texttt{PeerId} / \texttt{TopicId}
(types)}{PeerId / TopicId (types)}}\label{id-types}

In the public API, peer IDs and topic IDs are \textbf{66-char lowercase
hex strings}. The kernel holds them as \texttt{bigint} internally;
conversion happens at the API boundary. Pass strings unless a signature
explicitly says BigInt (\texttt{peer.lookup}).

\begin{center}\rule{0.5\linewidth}{0.5pt}\end{center}

\subsection{2. Identity}\label{identity}

Two factories. \texttt{createNodeIdentity} is the connection key
(location + nodeId, used for the handshake and routing).
\texttt{createAuthorIdentity} is the publish key (no location, used to
sign messages). They are separate on purpose (key separation): a peer
holds a node identity and signs each publish with an author identity
supplied per call.

\subsubsection{\texorpdfstring{\texttt{createNodeIdentity(\{\ lat,\ lng,\ extractable?\ \})}
→
\texttt{Promise\textless{}Identity\textgreater{}}}{createNodeIdentity(\{ lat, lng, extractable? \}) → Promise\textless Identity\textgreater{}}}\label{createnodeidentity-lat-lng-extractable-promiseidentity}

Generate a fresh node identity: a new Ed25519 keypair plus the 264-bit
nodeId derived from the public key + region.

{\def\LTcaptype{none} % do not increment counter
\begin{longtable}[]{@{}
  >{\raggedright\arraybackslash}p{(\linewidth - 4\tabcolsep) * \real{0.3333}}
  >{\raggedright\arraybackslash}p{(\linewidth - 4\tabcolsep) * \real{0.3333}}
  >{\raggedright\arraybackslash}p{(\linewidth - 4\tabcolsep) * \real{0.3333}}@{}}
\toprule\noalign{}
\begin{minipage}[b]{\linewidth}\raggedright
Param
\end{minipage} & \begin{minipage}[b]{\linewidth}\raggedright
Type
\end{minipage} & \begin{minipage}[b]{\linewidth}\raggedright
Notes
\end{minipage} \\
\midrule\noalign{}
\endhead
\bottomrule\noalign{}
\endlastfoot
\texttt{lat} & \texttt{number} & required --- latitude of the node's
region \\
\texttt{lng} & \texttt{number} & required --- longitude \\
\texttt{extractable} & \texttt{boolean} & default \texttt{true}. Pass
\texttt{false} for an ephemeral browser identity so the signing key
can't be exported (XSS can't exfiltrate it). Leave \texttt{true} if you
will \texttt{dumpIdentity} it. \\
\end{longtable}
}

\textbf{Returns} an \hyperref[identity-type]{\texttt{Identity}}. The top
byte of \texttt{.id} is the S2 cell for \texttt{(lat,\ lng)}.

\textbf{Throws} \texttt{IdentityError}:

\begin{itemize}
\tightlist
\item
  \texttt{IDENTITY\_INVALID\_FORMAT} --- \texttt{lat}/\texttt{lng} not
  numbers.
\item
  \texttt{IDENTITY\_KEYGEN\_FAILED} --- Web Crypto Ed25519
  \texttt{generateKey} failed.
\end{itemize}

\begin{Shaded}
\begin{Highlighting}[]
\ImportTok{import}\NormalTok{ \{ createNodeIdentity \} }\ImportTok{from} \StringTok{\textquotesingle{}@axona/protocol\textquotesingle{}}\OperatorTok{;}

\KeywordTok{const}\NormalTok{ node }\OperatorTok{=} \ControlFlowTok{await} \FunctionTok{createNodeIdentity}\NormalTok{(\{ }\DataTypeTok{lat}\OperatorTok{:} \FloatTok{38.0}\OperatorTok{,} \DataTypeTok{lng}\OperatorTok{:} \OperatorTok{{-}}\FloatTok{77.0}\NormalTok{ \})}\OperatorTok{;}
\BuiltInTok{console}\OperatorTok{.}\FunctionTok{log}\NormalTok{(node}\OperatorTok{.}\AttributeTok{id}\NormalTok{)}\OperatorTok{;}             \CommentTok{// \textquotesingle{}df…\textquotesingle{}  (66 hex chars)}
\BuiltInTok{console}\OperatorTok{.}\FunctionTok{log}\NormalTok{(node}\OperatorTok{.}\AttributeTok{id}\OperatorTok{.}\FunctionTok{slice}\NormalTok{(}\DecValTok{0}\OperatorTok{,} \DecValTok{2}\NormalTok{))}\OperatorTok{;} \CommentTok{// region byte}
\end{Highlighting}
\end{Shaded}

\subsubsection{\texorpdfstring{\texttt{createAuthorIdentity(\{\ persistAs?,\ store?,\ extractable?\ \})}
→
\texttt{Promise\textless{}AuthorIdentity\textgreater{}}}{createAuthorIdentity(\{ persistAs?, store?, extractable? \}) → Promise\textless AuthorIdentity\textgreater{}}}\label{createauthoridentity-persistas-store-extractable-promiseauthoridentity}

Mint (or load-or-create) the authorship key whose public key is the
Author ID. Durability is the only real choice: ephemeral (unlinkable) or
persisted (a recognizable author across sessions, able to retract).

{\def\LTcaptype{none} % do not increment counter
\begin{longtable}[]{@{}
  >{\raggedright\arraybackslash}p{(\linewidth - 4\tabcolsep) * \real{0.3333}}
  >{\raggedright\arraybackslash}p{(\linewidth - 4\tabcolsep) * \real{0.3333}}
  >{\raggedright\arraybackslash}p{(\linewidth - 4\tabcolsep) * \real{0.3333}}@{}}
\toprule\noalign{}
\begin{minipage}[b]{\linewidth}\raggedright
Param
\end{minipage} & \begin{minipage}[b]{\linewidth}\raggedright
Type
\end{minipage} & \begin{minipage}[b]{\linewidth}\raggedright
Notes
\end{minipage} \\
\midrule\noalign{}
\endhead
\bottomrule\noalign{}
\endlastfoot
\texttt{persistAs} & \texttt{string} & storage key; load-or-create then
save. Omit for an ephemeral author. \\
\texttt{store} & \texttt{\{\ get,\ set\ \}} & custom store; defaults to
browser \texttt{localStorage} when \texttt{persistAs} is set. \\
\texttt{extractable} & \texttt{boolean} & default \texttt{true};
\textbf{forced \texttt{true}} when \texttt{persistAs} is set (a
persisted key must be exportable). \\
\end{longtable}
}

\textbf{Returns} an
\hyperref[author-identity-type]{\texttt{AuthorIdentity}}.

\textbf{Throws} \texttt{IdentityError}
(\texttt{IDENTITY\_KEYGEN\_FAILED}, \texttt{IDENTITY\_LOAD\_FAILED},
\texttt{IDENTITY\_INVALID\_FORMAT}).

\begin{Shaded}
\begin{Highlighting}[]
\ImportTok{import}\NormalTok{ \{ createAuthorIdentity \} }\ImportTok{from} \StringTok{\textquotesingle{}@axona/protocol\textquotesingle{}}\OperatorTok{;}

\KeywordTok{const}\NormalTok{ ephemeral }\OperatorTok{=} \ControlFlowTok{await} \FunctionTok{createAuthorIdentity}\NormalTok{()}\OperatorTok{;}              \CommentTok{// unlinkable, one session}
\KeywordTok{const}\NormalTok{ durable   }\OperatorTok{=} \ControlFlowTok{await} \FunctionTok{createAuthorIdentity}\NormalTok{(\{ }\DataTypeTok{persistAs}\OperatorTok{:} \StringTok{\textquotesingle{}me\textquotesingle{}}\NormalTok{ \})}\OperatorTok{;} \CommentTok{// recognizable across reloads}
\BuiltInTok{console}\OperatorTok{.}\FunctionTok{log}\NormalTok{(durable}\OperatorTok{.}\AttributeTok{authorId}\NormalTok{)}\OperatorTok{;}   \CommentTok{// 64{-}hex Author ID — share this as \textasciigrave{}owner\textasciigrave{}}
\end{Highlighting}
\end{Shaded}

\subsubsection{\texorpdfstring{\texttt{dumpIdentity(identity)} →
\texttt{Promise\textless{}IdentityEnvelope\textgreater{}}
\emph{(advanced)}}{dumpIdentity(identity) → Promise\textless IdentityEnvelope\textgreater{} (advanced)}}\label{dumpidentityidentity-promiseidentityenvelope-advanced}

Serialize a \textbf{node} identity to a JSON-safe envelope (the
\texttt{privateKey} is exported as base64 PKCS\#8). Loses the in-memory
\texttt{CryptoKey} handle; reconstruct with \texttt{loadIdentity}.

\begin{Shaded}
\begin{Highlighting}[]
\KeywordTok{const}\NormalTok{ env }\OperatorTok{=} \ControlFlowTok{await} \FunctionTok{dumpIdentity}\NormalTok{(node)}\OperatorTok{;}
\CommentTok{// \{ id, pubkey, privkey, region: \{ lat, lng \}, createdAt, pow \}}
\NormalTok{localStorage}\OperatorTok{.}\FunctionTok{setItem}\NormalTok{(}\StringTok{\textquotesingle{}node{-}identity\textquotesingle{}}\OperatorTok{,} \BuiltInTok{JSON}\OperatorTok{.}\FunctionTok{stringify}\NormalTok{(env))}\OperatorTok{;}
\end{Highlighting}
\end{Shaded}

\textbf{Throws} \texttt{IdentityError(IDENTITY\_LOAD\_FAILED)} if the
key can't be exported.

\subsubsection{\texorpdfstring{\texttt{loadIdentity(envelope)} →
\texttt{Promise\textless{}Identity\textgreater{}}
\emph{(advanced)}}{loadIdentity(envelope) → Promise\textless Identity\textgreater{} (advanced)}}\label{loadidentityenvelope-promiseidentity-advanced}

Inverse of \texttt{dumpIdentity}. Reconstructs the full node
\texttt{Identity}, verifying that (a) the stored id matches the
freshly-derived id and (b) the private key corresponds to the public key
(a sign→verify probe).

\begin{Shaded}
\begin{Highlighting}[]
\KeywordTok{const}\NormalTok{ env }\OperatorTok{=} \BuiltInTok{JSON}\OperatorTok{.}\FunctionTok{parse}\NormalTok{(localStorage}\OperatorTok{.}\FunctionTok{getItem}\NormalTok{(}\StringTok{\textquotesingle{}node{-}identity\textquotesingle{}}\NormalTok{))}\OperatorTok{;}
\KeywordTok{const}\NormalTok{ node }\OperatorTok{=} \ControlFlowTok{await} \FunctionTok{loadIdentity}\NormalTok{(env)}\OperatorTok{;}
\end{Highlighting}
\end{Shaded}

\textbf{Throws} \texttt{IdentityError}:

\begin{itemize}
\tightlist
\item
  \texttt{IDENTITY\_INVALID\_FORMAT} --- malformed envelope, id
  mismatch, or private/public key mismatch.
\item
  \texttt{IDENTITY\_LOAD\_FAILED} --- PKCS\#8 import failed.
\end{itemize}

\begin{quote}
Author-identity persistence is handled inside
\texttt{createAuthorIdentity(\{\ persistAs\ \})} --- you don't dump/load
author identities by hand. \texttt{dumpIdentity}/\texttt{loadIdentity}
are for node identities (or when you manage your own storage layer).
\end{quote}

\subsubsection{\texorpdfstring{\texttt{computeNodeId(pubkeyBytes,\ lat,\ lng)}
→ \texttt{Promise\textless{}string\textgreater{}}
\emph{(advanced)}}{computeNodeId(pubkeyBytes, lat, lng) → Promise\textless string\textgreater{} (advanced)}}\label{computenodeidpubkeybytes-lat-lng-promisestring-advanced}

\subsubsection{\texorpdfstring{\texttt{computeNodeIdBigInt(pubkeyBytes,\ lat,\ lng)}
→ \texttt{Promise\textless{}bigint\textgreater{}}
\emph{(advanced)}}{computeNodeIdBigInt(pubkeyBytes, lat, lng) → Promise\textless bigint\textgreater{} (advanced)}}\label{computenodeidbigintpubkeybytes-lat-lng-promisebigint-advanced}

Derive (or verify) a 264-bit nodeId from a raw public key + region.
\texttt{computeNodeId} returns the 66-hex form;
\texttt{computeNodeIdBigInt} returns the BigInt. Used to validate a
claimed nodeId against a pubkey + region without minting a fresh
identity.

\begin{Shaded}
\begin{Highlighting}[]
\KeywordTok{const}\NormalTok{ id }\OperatorTok{=} \ControlFlowTok{await} \FunctionTok{computeNodeId}\NormalTok{(node}\OperatorTok{.}\AttributeTok{pubkey}\OperatorTok{,} \FloatTok{38.0}\OperatorTok{,} \OperatorTok{{-}}\FloatTok{77.0}\NormalTok{)}\OperatorTok{;}
\BuiltInTok{console}\OperatorTok{.}\FunctionTok{log}\NormalTok{(id }\OperatorTok{===}\NormalTok{ node}\OperatorTok{.}\AttributeTok{id}\NormalTok{)}\OperatorTok{;}   \CommentTok{// true}
\end{Highlighting}
\end{Shaded}

\begin{center}\rule{0.5\linewidth}{0.5pt}\end{center}

\subsection{3. AxonaPeer construction +
lifecycle}\label{axonapeer-construction-lifecycle}

\subsubsection{\texorpdfstring{\texttt{connect(\{\ bridge,\ location,\ author?,\ k?,\ ready?,\ transport?,\ nodeIdentity?,\ web?\ \})}
→
\texttt{Promise\textless{}\{\ peer,\ author,\ nodeIdentity,\ transport,\ status,\ disconnect\ \}\textgreater{}}
\emph{(kernel
4.16.0)}}{connect(\{ bridge, location, author?, k?, ready?, transport?, nodeIdentity?, web? \}) → Promise\textless\{ peer, author, nodeIdentity, transport, status, disconnect \}\textgreater{} (kernel 4.16.0)}}\label{connect-bridge-location-author-k-ready-transport-nodeidentity-web-promise-peer-author-nodeidentity-transport-status-disconnect-kernel-4.16.0}

The one-call bootstrap --- import from the \texttt{connect.js} subpath:

\begin{Shaded}
\begin{Highlighting}[]
\ImportTok{import}\NormalTok{ \{ connect \} }\ImportTok{from} \StringTok{\textquotesingle{}@axona/protocol/connect.js\textquotesingle{}}\OperatorTok{;}
\KeywordTok{const}\NormalTok{ \{ peer}\OperatorTok{,}\NormalTok{ author}\OperatorTok{,}\NormalTok{ status}\OperatorTok{,}\NormalTok{ disconnect \} }\OperatorTok{=} \ControlFlowTok{await} \FunctionTok{connect}\NormalTok{(\{}
  \DataTypeTok{bridge}\OperatorTok{:} \StringTok{\textquotesingle{}wss://testnet.axona.net\textquotesingle{}}\OperatorTok{,} \DataTypeTok{location}\OperatorTok{:}\NormalTok{ \{ }\DataTypeTok{lat}\OperatorTok{:} \DecValTok{38}\OperatorTok{,} \DataTypeTok{lng}\OperatorTok{:} \OperatorTok{{-}}\DecValTok{77}\NormalTok{ \} \})}\OperatorTok{;}
\end{Highlighting}
\end{Shaded}

Mints the connection identity (and, by default, an ephemeral author),
builds the web transport + \texttt{NeuronNode} + \texttt{AxonaDomain} +
\texttt{AxonaPeer}, starts everything, and awaits \texttt{peer.ready()}.
Pure sugar over the constructors below --- use them directly for custom
wiring.

{\def\LTcaptype{none} % do not increment counter
\begin{longtable}[]{@{}
  >{\raggedright\arraybackslash}p{(\linewidth - 4\tabcolsep) * \real{0.3333}}
  >{\raggedright\arraybackslash}p{(\linewidth - 4\tabcolsep) * \real{0.3333}}
  >{\raggedright\arraybackslash}p{(\linewidth - 4\tabcolsep) * \real{0.3333}}@{}}
\toprule\noalign{}
\begin{minipage}[b]{\linewidth}\raggedright
Param
\end{minipage} & \begin{minipage}[b]{\linewidth}\raggedright
Type
\end{minipage} & \begin{minipage}[b]{\linewidth}\raggedright
Notes
\end{minipage} \\
\midrule\noalign{}
\endhead
\bottomrule\noalign{}
\endlastfoot
\texttt{bridge} & string & \texttt{wss://} bridge URL. Required unless
\texttt{transport} injected. \\
\texttt{location} & \texttt{\{\ lat,\ lng\ \}} & The node's real
location. Required unless \texttt{nodeIdentity} injected. \\
\texttt{author} & \texttt{true} \textbar{} string \textbar{} identity
\textbar{} \texttt{false} & \texttt{true} (default) ephemeral author; a
string → durable via \texttt{persistAs}; an author identity → used
as-is; \texttt{false} → none. \\
\texttt{k} & number (20) & Routing closest-set size. \\
\texttt{ready} & object \textbar{} \texttt{false} & Forwarded to
\texttt{peer.ready()}; \texttt{false} skips the wait (\texttt{status} is
\texttt{null}). \\
\texttt{transport} / \texttt{nodeIdentity} & --- & Injection for tests,
sim, custom stacks; the web transport is only loaded when no
\texttt{transport} is given. \\
\texttt{web} & object & Extra options forwarded to the
\texttt{webTransport} factory. \\
\end{longtable}
}

\texttt{disconnect()} performs \texttt{peer.leave()} +
\texttt{peer.stop()} + \texttt{transport.stop()}, best-effort.
\texttt{connect} lives outside the main barrel (the barrel stays
environment-neutral); the import path is
\texttt{@axona/protocol/connect.js}.

\texttt{AxonaPeer} is the per-node DHT contract implementation --- one
instance per running node.

\subsubsection{\texorpdfstring{\texttt{new\ AxonaPeer(\{\ nodeIdentity,\ domain,\ transport,\ ...\ \})}}{new AxonaPeer(\{ nodeIdentity, domain, transport, ... \})}}\label{new-axonapeer-nodeidentity-domain-transport-...}

{\def\LTcaptype{none} % do not increment counter
\begin{longtable}[]{@{}
  >{\raggedright\arraybackslash}p{(\linewidth - 4\tabcolsep) * \real{0.3333}}
  >{\raggedright\arraybackslash}p{(\linewidth - 4\tabcolsep) * \real{0.3333}}
  >{\raggedright\arraybackslash}p{(\linewidth - 4\tabcolsep) * \real{0.3333}}@{}}
\toprule\noalign{}
\begin{minipage}[b]{\linewidth}\raggedright
Param
\end{minipage} & \begin{minipage}[b]{\linewidth}\raggedright
Type
\end{minipage} & \begin{minipage}[b]{\linewidth}\raggedright
Notes
\end{minipage} \\
\midrule\noalign{}
\endhead
\bottomrule\noalign{}
\endlastfoot
\texttt{node} & \texttt{NeuronNode} & \textbf{required} --- the per-node
state object (\texttt{\{\ id,\ alive,\ transport,\ synaptome\ \}}). \\
\texttt{nodeIdentity} & \texttt{Identity} & the connection key (from
\texttt{createNodeIdentity}). Used for the handshake, routing, and
signing the node's own actions. It \textbf{never} signs a publish. \\
\texttt{transport} & \texttt{Transport} & the transport instance
(web/sim/node). The peer uses it for
\texttt{send}/\texttt{notify}/\texttt{lookup}. \\
\texttt{domain} & \texttt{AxonaDomain} & shared routing parameters.
\textbf{\texttt{domain} or \texttt{engine} is required.} \\
\texttt{engine} & object & legacy simulator engine; pass instead of
\texttt{domain} only in sim/tests. \\
\texttt{axonaManager} & \texttt{AxonaManager} & pre-built pub/sub
manager; if omitted, resolved on first \texttt{pub}/\texttt{sub}. \\
\texttt{persist} & \texttt{PersistenceAdapter} & enables
auto-checkpointing of identity, subscriptions, synaptome, and hosting
state. \\
\texttt{maxPublishBytes} & \texttt{number} & per-publish cap; clamped to
the WebRTC-interop floor (16 KiB) and never above the absolute ingress
cap. Override only for known-homogeneous fleets. \\
\texttt{synaptomeMaintain} &
\texttt{boolean\ \textbackslash{}\textbar{}\ \{\ kNear?,\ intervalMs?,\ maxPerTick?\ \}}
& \emph{(advanced; v4.9.0)} opt into continuous \textbf{near-quota
maintenance} --- the peer keeps its \texttt{kNear} (≈5) XOR-nearest
``successor'' links filled through churn so greedy routing's last mile
stays complete. \texttt{true} uses the defaults
\texttt{\{\ kNear:\ 5,\ intervalMs:\ 15000,\ maxPerTick:\ 3\ \}}. Off by
default; the standard \texttt{axona-peer}/relay builds set it.
Long-range ``fingers'' are handled separately by the routing anneal. \\
\end{longtable}
}

The production path is
\texttt{\{\ nodeIdentity,\ domain,\ transport\ \}}. The
\texttt{\{\ engine\ \}} / \texttt{\{\ axonaManager\ \}} forms are for
the simulator and tests.

\begin{Shaded}
\begin{Highlighting}[]
\ImportTok{import}\NormalTok{ \{ AxonaPeer}\OperatorTok{,}\NormalTok{ AxonaDomain}\OperatorTok{,}\NormalTok{ NeuronNode}\OperatorTok{,}\NormalTok{ createNodeIdentity \} }\ImportTok{from} \StringTok{\textquotesingle{}@axona/protocol\textquotesingle{}}\OperatorTok{;}

\KeywordTok{const}\NormalTok{ node }\OperatorTok{=} \ControlFlowTok{await} \FunctionTok{createNodeIdentity}\NormalTok{(\{ }\DataTypeTok{lat}\OperatorTok{:} \DecValTok{38}\OperatorTok{,} \DataTypeTok{lng}\OperatorTok{:} \OperatorTok{{-}}\DecValTok{77}\NormalTok{ \})}\OperatorTok{;}
\KeywordTok{const}\NormalTok{ peer }\OperatorTok{=} \KeywordTok{new} \FunctionTok{AxonaPeer}\NormalTok{(\{}
  \DataTypeTok{nodeIdentity}\OperatorTok{:}\NormalTok{ node}\OperatorTok{,}
  \DataTypeTok{domain}\OperatorTok{:}       \KeywordTok{new} \FunctionTok{AxonaDomain}\NormalTok{(\{ }\DataTypeTok{k}\OperatorTok{:} \DecValTok{20}\NormalTok{ \})}\OperatorTok{,}
  \DataTypeTok{node}\OperatorTok{:}         \KeywordTok{new} \FunctionTok{NeuronNode}\NormalTok{(\{ }\DataTypeTok{id}\OperatorTok{:} \BuiltInTok{BigInt}\NormalTok{(}\StringTok{\textquotesingle{}0x\textquotesingle{}} \OperatorTok{+}\NormalTok{ node}\OperatorTok{.}\AttributeTok{id}\NormalTok{)}\OperatorTok{,} \DataTypeTok{lat}\OperatorTok{:} \DecValTok{38}\OperatorTok{,} \DataTypeTok{lng}\OperatorTok{:} \OperatorTok{{-}}\DecValTok{77}\NormalTok{ \})}\OperatorTok{,}
\NormalTok{  transport}\OperatorTok{,}
\NormalTok{\})}\OperatorTok{;}
\end{Highlighting}
\end{Shaded}

\textbf{Throws} plain \texttt{Error} if \texttt{node} is missing, or if
neither \texttt{engine} nor \texttt{domain} is supplied.

\subsubsection{\texorpdfstring{\texttt{peer.start()} →
\texttt{Promise\textless{}void\textgreater{}}}{peer.start() → Promise\textless void\textgreater{}}}\label{peer.start-promisevoid}

Bring the peer up --- installs wire-level routing handlers, the delivery
hook, the peer-bound/peer-died eviction handlers, and (if
\texttt{persist} is wired) hydrates persisted state. Idempotent.

\subsubsection{\texorpdfstring{\texttt{peer.ready(\{\ minPeers?,\ timeoutMs?,\ stableMs?,\ pollMs?\ \})}
→
\texttt{Promise\textless{}\{\ ready,\ peers,\ ms,\ reason\ \}\textgreater{}}}{peer.ready(\{ minPeers?, timeoutMs?, stableMs?, pollMs? \}) → Promise\textless\{ ready, peers, ms, reason \}\textgreater{}}}\label{peer.ready-minpeers-timeoutms-stablems-pollms-promise-ready-peers-ms-reason}

Wait for the routing mesh to warm up before you
\texttt{pub}/\texttt{sub}. \texttt{start()} returns as soon as local
state is set up, but mesh handshakes complete asynchronously over the
next few seconds; publishing/subscribing on a cold synaptome routes to
too few root axons. \texttt{ready()} resolves once the synaptome reaches
\texttt{minPeers}, or --- failing that --- settles at a \textbf{stable
non-zero plateau} (unchanged for \texttt{stableMs}), or the
\texttt{timeoutMs} budget elapses. Use it instead of hand-rolling a
\texttt{peer.peers().length} / \texttt{node.synaptome.size} poll.

{\def\LTcaptype{none} % do not increment counter
\begin{longtable}[]{@{}
  >{\raggedright\arraybackslash}p{(\linewidth - 6\tabcolsep) * \real{0.2500}}
  >{\raggedright\arraybackslash}p{(\linewidth - 6\tabcolsep) * \real{0.2500}}
  >{\raggedright\arraybackslash}p{(\linewidth - 6\tabcolsep) * \real{0.2500}}
  >{\raggedright\arraybackslash}p{(\linewidth - 6\tabcolsep) * \real{0.2500}}@{}}
\toprule\noalign{}
\begin{minipage}[b]{\linewidth}\raggedright
Param
\end{minipage} & \begin{minipage}[b]{\linewidth}\raggedright
Type
\end{minipage} & \begin{minipage}[b]{\linewidth}\raggedright
Default
\end{minipage} & \begin{minipage}[b]{\linewidth}\raggedright
Notes
\end{minipage} \\
\midrule\noalign{}
\endhead
\bottomrule\noalign{}
\endlastfoot
\texttt{minPeers} & \texttt{number} & \texttt{4} & Resolve
\texttt{ready:true} as soon as the synaptome reaches this size. \\
\texttt{timeoutMs} & \texttt{number} & \texttt{10000} & Upper bound; on
expiry resolves with \texttt{ready} = (peers \textgreater{} 0). \\
\texttt{stableMs} & \texttt{number} & \texttt{1500} & If below
\texttt{minPeers} but non-zero and unchanged this long, resolve
\texttt{ready:true}
(\texttt{reason:\textquotesingle{}stable\textquotesingle{}}). \\
\texttt{pollMs} & \texttt{number} & \texttt{150} & Synaptome poll
interval. \\
\end{longtable}
}

Returns
\texttt{\{\ ready:boolean,\ peers:number,\ ms:number,\ reason:\textquotesingle{}minPeers\textquotesingle{}\textbar{}\textquotesingle{}stable\textquotesingle{}\textbar{}\textquotesingle{}timeout\textquotesingle{}\ \}}.
Never rejects --- a cold/isolated start resolves
\texttt{\{\ ready:false,\ …,\ reason:\textquotesingle{}timeout\textquotesingle{}\ \}}
so you can surface ``still connecting'' rather than throw.

\begin{Shaded}
\begin{Highlighting}[]
\ControlFlowTok{await}\NormalTok{ peer}\OperatorTok{.}\FunctionTok{start}\NormalTok{()}\OperatorTok{;}
\KeywordTok{const}\NormalTok{ status }\OperatorTok{=} \ControlFlowTok{await}\NormalTok{ peer}\OperatorTok{.}\FunctionTok{ready}\NormalTok{()}\OperatorTok{;}          \CommentTok{// defaults: 4 peers or a stable plateau, ≤10s}
\ControlFlowTok{if}\NormalTok{ (}\OperatorTok{!}\NormalTok{status}\OperatorTok{.}\AttributeTok{ready}\NormalTok{) }\BuiltInTok{console}\OperatorTok{.}\FunctionTok{warn}\NormalTok{(}\StringTok{\textquotesingle{}mesh still warming — publishes may be thin\textquotesingle{}}\NormalTok{)}\OperatorTok{;}
\end{Highlighting}
\end{Shaded}

\subsubsection{\texorpdfstring{\texttt{peer.stop()} →
\texttt{Promise\textless{}void\textgreater{}}}{peer.stop() → Promise\textless void\textgreater{}}}\label{peer.stop-promisevoid}

Tear down: release event listeners and the peer-lifecycle hooks.
Idempotent. (For a graceful network exit, prefer \texttt{peer.leave}.)

\subsubsection{\texorpdfstring{\texttt{peer.join(sponsor?)} →
\texttt{Promise\textless{}void\textgreater{}}}{peer.join(sponsor?) → Promise\textless void\textgreater{}}}\label{peer.joinsponsor-promisevoid}

Bootstrap into the mesh. Calls \texttt{start()} first if needed, brings
up the transport, and --- if a \texttt{sponsor} is given --- opens a
channel to it and seeds the synaptome.

{\def\LTcaptype{none} % do not increment counter
\begin{longtable}[]{@{}
  >{\raggedright\arraybackslash}p{(\linewidth - 4\tabcolsep) * \real{0.3333}}
  >{\raggedright\arraybackslash}p{(\linewidth - 4\tabcolsep) * \real{0.3333}}
  >{\raggedright\arraybackslash}p{(\linewidth - 4\tabcolsep) * \real{0.3333}}@{}}
\toprule\noalign{}
\begin{minipage}[b]{\linewidth}\raggedright
Param
\end{minipage} & \begin{minipage}[b]{\linewidth}\raggedright
Type
\end{minipage} & \begin{minipage}[b]{\linewidth}\raggedright
Notes
\end{minipage} \\
\midrule\noalign{}
\endhead
\bottomrule\noalign{}
\endlastfoot
\texttt{sponsor} & \texttt{string} & 66-char hex nodeId of a known peer.
Omit to start standalone (inbound connections will populate the
synaptome). \\
\end{longtable}
}

\begin{Shaded}
\begin{Highlighting}[]
\ControlFlowTok{await}\NormalTok{ peer}\OperatorTok{.}\FunctionTok{join}\NormalTok{()}\OperatorTok{;}                 \CommentTok{// standalone; wait for inbound peers}
\ControlFlowTok{await}\NormalTok{ peer}\OperatorTok{.}\FunctionTok{join}\NormalTok{(bridgeNodeIdHex)}\OperatorTok{;}  \CommentTok{// open a channel to a sponsor and seed}
\end{Highlighting}
\end{Shaded}

\textbf{Throws} \texttt{TransportError}:

\begin{itemize}
\tightlist
\item
  \texttt{TRANSPORT\_NOT\_STARTED} --- \texttt{transport.start} failed,
  or \texttt{sponsor} is not 66-char hex.
\item
  \texttt{TRANSPORT\_PEER\_UNREACHABLE} --- the sponsor isn't reachable.
\end{itemize}

\subsubsection{\texorpdfstring{\texttt{peer.leave(\{\ drain?,\ notify?,\ timeoutMs?\ \})}
→
\texttt{Promise\textless{}void\textgreater{}}}{peer.leave(\{ drain?, notify?, timeoutMs? \}) → Promise\textless void\textgreater{}}}\label{peer.leave-drain-notify-timeoutms-promisevoid}

Graceful shutdown.

{\def\LTcaptype{none} % do not increment counter
\begin{longtable}[]{@{}
  >{\raggedright\arraybackslash}p{(\linewidth - 6\tabcolsep) * \real{0.2500}}
  >{\raggedright\arraybackslash}p{(\linewidth - 6\tabcolsep) * \real{0.2500}}
  >{\raggedright\arraybackslash}p{(\linewidth - 6\tabcolsep) * \real{0.2500}}
  >{\raggedright\arraybackslash}p{(\linewidth - 6\tabcolsep) * \real{0.2500}}@{}}
\toprule\noalign{}
\begin{minipage}[b]{\linewidth}\raggedright
Param
\end{minipage} & \begin{minipage}[b]{\linewidth}\raggedright
Type
\end{minipage} & \begin{minipage}[b]{\linewidth}\raggedright
Default
\end{minipage} & \begin{minipage}[b]{\linewidth}\raggedright
Notes
\end{minipage} \\
\midrule\noalign{}
\endhead
\bottomrule\noalign{}
\endlastfoot
\texttt{notify} & \texttt{boolean} & \texttt{true} & send a
\texttt{peer-leaving} notification to every synaptome peer so they drop
us proactively (sub-second handoff instead of heartbeat timeout). \\
\texttt{drain} & \texttt{boolean} & \texttt{true} & pause briefly for
in-flight publishes/pulls to settle before closing. \\
\texttt{timeoutMs} & \texttt{number} & \texttt{5000} & bound on the
drain wait. \\
\end{longtable}
}

Closes the transport last, force-flushes persistence, and stops
listeners.

\begin{Shaded}
\begin{Highlighting}[]
\ControlFlowTok{await}\NormalTok{ peer}\OperatorTok{.}\FunctionTok{leave}\NormalTok{(\{ }\DataTypeTok{drain}\OperatorTok{:} \KeywordTok{true}\OperatorTok{,} \DataTypeTok{notify}\OperatorTok{:} \KeywordTok{true}\OperatorTok{,} \DataTypeTok{timeoutMs}\OperatorTok{:} \DecValTok{3000}\NormalTok{ \})}\OperatorTok{;}
\end{Highlighting}
\end{Shaded}

\subsubsection{\texorpdfstring{\texttt{peer.getNodeId()} →
\texttt{string\ \textbar{}\ bigint}}{peer.getNodeId() → string \textbar{} bigint}}\label{peer.getnodeid-string-bigint}

Returns whatever was passed as \texttt{node.id} at construction time
(BigInt if you followed the convention).

\subsubsection{\texorpdfstring{\texttt{peer.snapshot()} →
\texttt{Promise\textless{}object\textgreater{}}}{peer.snapshot() → Promise\textless object\textgreater{}}}\label{peer.snapshot-promiseobject}

Serialize this peer's state to a JSON-safe envelope:
\texttt{\{\ formatVersion:\ \textquotesingle{}1.0\textquotesingle{},\ snapshotAt,\ wireVersion,\ identity,\ synaptome,\ subscriptions\ \}}.
Store it however you like (encrypted, synced, custom DB).

\begin{Shaded}
\begin{Highlighting}[]
\KeywordTok{const}\NormalTok{ state }\OperatorTok{=} \ControlFlowTok{await}\NormalTok{ peer}\OperatorTok{.}\FunctionTok{snapshot}\NormalTok{()}\OperatorTok{;}
\NormalTok{localStorage}\OperatorTok{.}\FunctionTok{setItem}\NormalTok{(}\StringTok{\textquotesingle{}peer{-}state\textquotesingle{}}\OperatorTok{,} \BuiltInTok{JSON}\OperatorTok{.}\FunctionTok{stringify}\NormalTok{(state))}\OperatorTok{;}
\end{Highlighting}
\end{Shaded}

\subsubsection{\texorpdfstring{\texttt{AxonaPeer.fromSnapshot(state,\ opts)}
→ \texttt{Promise\textless{}AxonaPeer\textgreater{}}
\emph{(static)}}{AxonaPeer.fromSnapshot(state, opts) → Promise\textless AxonaPeer\textgreater{} (static)}}\label{axonapeer.fromsnapshotstate-opts-promiseaxonapeer-static}

Reconstruct a peer from a \texttt{snapshot()} envelope. The returned
peer is constructed with identity / synaptome pre-loaded, but the
transport is \textbf{not} started --- call \texttt{peer.join(sponsor?)}
to bring it up. Subscription \textbf{handlers are not restored}
(functions don't serialize); the restored list is exposed at
\texttt{peer.pendingSubscriptions} so you can re-register them with
\texttt{peer.sub}.

\begin{Shaded}
\begin{Highlighting}[]
\KeywordTok{const}\NormalTok{ state }\OperatorTok{=} \BuiltInTok{JSON}\OperatorTok{.}\FunctionTok{parse}\NormalTok{(localStorage}\OperatorTok{.}\FunctionTok{getItem}\NormalTok{(}\StringTok{\textquotesingle{}peer{-}state\textquotesingle{}}\NormalTok{))}\OperatorTok{;}
\KeywordTok{const}\NormalTok{ peer  }\OperatorTok{=} \ControlFlowTok{await}\NormalTok{ AxonaPeer}\OperatorTok{.}\FunctionTok{fromSnapshot}\NormalTok{(state}\OperatorTok{,}\NormalTok{ \{ domain}\OperatorTok{,}\NormalTok{ node}\OperatorTok{,}\NormalTok{ transport \})}\OperatorTok{;}
\ControlFlowTok{for}\NormalTok{ (}\KeywordTok{const}\NormalTok{ s }\KeywordTok{of}\NormalTok{ peer}\OperatorTok{.}\AttributeTok{pendingSubscriptions}\NormalTok{) \{}
  \ControlFlowTok{await}\NormalTok{ peer}\OperatorTok{.}\FunctionTok{sub}\NormalTok{(s}\OperatorTok{.}\AttributeTok{topic}\OperatorTok{,}\NormalTok{ onMsg}\OperatorTok{,}\NormalTok{ \{ }\DataTypeTok{since}\OperatorTok{:}\NormalTok{ s}\OperatorTok{.}\AttributeTok{since}\NormalTok{ \})}\OperatorTok{;}
\NormalTok{\}}
\ControlFlowTok{await}\NormalTok{ peer}\OperatorTok{.}\FunctionTok{join}\NormalTok{()}\OperatorTok{;}
\end{Highlighting}
\end{Shaded}

\textbf{Throws} \texttt{TypeError} / \texttt{RangeError} if
\texttt{state} isn't a \texttt{1.0}-format snapshot object.

\begin{center}\rule{0.5\linewidth}{0.5pt}\end{center}

\subsection{4. Pub/sub}\label{pubsub}

\subsubsection{4.1 Topic addressing}\label{topic-addressing}

\emph{(Folded in from the former standalone ``Topic IDs'' note.)}

A topic is not a bare string --- it's a small \textbf{descriptor}, and
the topic ID is a hash of that descriptor with a one-byte region prefix:

\begin{verbatim}
topicId = regionByte || SHA-256(canonical({ owner, name, write }))   // 33 bytes = 66 hex chars
\end{verbatim}

{\def\LTcaptype{none} % do not increment counter
\begin{longtable}[]{@{}
  >{\raggedright\arraybackslash}p{(\linewidth - 2\tabcolsep) * \real{0.5000}}
  >{\raggedright\arraybackslash}p{(\linewidth - 2\tabcolsep) * \real{0.5000}}@{}}
\toprule\noalign{}
\begin{minipage}[b]{\linewidth}\raggedright
Field
\end{minipage} & \begin{minipage}[b]{\linewidth}\raggedright
Meaning
\end{minipage} \\
\midrule\noalign{}
\endhead
\bottomrule\noalign{}
\endlastfoot
\texttt{region} & a real geographic cell; becomes the leading byte of
the ID \\
\texttt{owner} & an \textbf{Author ID} (a publish key), or absent \\
\texttt{name} & the human label --- \texttt{"lobby"},
\texttt{"profile"} \\
\texttt{write} & \texttt{"open"} or \texttt{"owner"} --- defaults by
whether \texttt{owner} is set \\
\end{longtable}
}

Because the ID is a pure function of those fields, anyone who knows the
fields computes the identical ID --- no registry, no coordination.

\textbf{The \texttt{region} field accepts a name, hex string, decimal
string, or number} --- all normalize to the same region byte:

\begin{Shaded}
\begin{Highlighting}[]
\NormalTok{\{ }\DataTypeTok{region}\OperatorTok{:} \StringTok{\textquotesingle{}useast\textquotesingle{}}\OperatorTok{,} \DataTypeTok{name}\OperatorTok{:} \StringTok{\textquotesingle{}lobby\textquotesingle{}}\NormalTok{ \}   }\CommentTok{// name}
\NormalTok{\{ }\DataTypeTok{region}\OperatorTok{:} \StringTok{\textquotesingle{}0x89\textquotesingle{}}\OperatorTok{,}   \DataTypeTok{name}\OperatorTok{:} \StringTok{\textquotesingle{}lobby\textquotesingle{}}\NormalTok{ \}   }\CommentTok{// hex string   {-}{-}\textbackslash{}  same topic ID}
\NormalTok{\{ }\DataTypeTok{region}\OperatorTok{:} \DecValTok{137}\OperatorTok{,}      \DataTypeTok{name}\OperatorTok{:} \StringTok{\textquotesingle{}lobby\textquotesingle{}}\NormalTok{ \}   }\CommentTok{// number       {-}{-}/  (leading byte 0x89)}
\end{Highlighting}
\end{Shaded}

Omit \texttt{region} and it defaults to the publisher's own node region
(top byte of its node ID). Unknown regions throw \texttt{RangeError}.

\textbf{\texttt{write} defaults by whether there's an \texttt{owner}:}

{\def\LTcaptype{none} % do not increment counter
\begin{longtable}[]{@{}
  >{\raggedright\arraybackslash}p{(\linewidth - 4\tabcolsep) * \real{0.3333}}
  >{\raggedright\arraybackslash}p{(\linewidth - 4\tabcolsep) * \real{0.3333}}
  >{\raggedright\arraybackslash}p{(\linewidth - 4\tabcolsep) * \real{0.3333}}@{}}
\toprule\noalign{}
\begin{minipage}[b]{\linewidth}\raggedright
Descriptor
\end{minipage} & \begin{minipage}[b]{\linewidth}\raggedright
Resolved \texttt{write}
\end{minipage} & \begin{minipage}[b]{\linewidth}\raggedright
Meaning
\end{minipage} \\
\midrule\noalign{}
\endhead
\bottomrule\noalign{}
\endlastfoot
\texttt{\{\ region,\ name\ \}} (no owner) & \texttt{open} & open lobby
--- anyone publishes \\
\texttt{\{\ region,\ owner,\ name\ \}} (no write) &
\textbf{\texttt{owner}} & owned --- only the owner publishes \\
\texttt{\{\ region,\ owner,\ name,\ write:\textquotesingle{}owner\textquotesingle{}\ \}}
& \texttt{owner} & same ID as the row above \\
\texttt{\{\ region,\ owner,\ name,\ write:\textquotesingle{}open\textquotesingle{}\ \}}
& \texttt{open} & owner-namespaced, anyone publishes (inbox/wall) \\
\texttt{\{\ region,\ name,\ write:\textquotesingle{}owner\textquotesingle{}\ \}}
(no owner) & \texttt{open} & no owner -\textgreater{} \texttt{write}
ignored \\
\end{longtable}
}

Two rules: \textbf{no \texttt{owner} -\textgreater{} the topic is open
and \texttt{write} is ignored}; \textbf{an \texttt{owner}
-\textgreater{} \texttt{write} defaults to
\texttt{\textquotesingle{}owner\textquotesingle{}}} (the safe default
--- forgetting \texttt{write} can never silently make an owned feed
world-writable).

\textbf{Share the ID for reading; share the descriptor for writing.}
\texttt{sub}, \texttt{pull}, and \texttt{metrics} accept the 66-hex ID
or a descriptor. \texttt{pub} (and the owner op \texttt{kill}) requires
the descriptor --- a bare ID is rejected, because a hash can't reveal
its \texttt{owner}, so the ID alone can't prove write authorization.

\paragraph{\texorpdfstring{\texttt{deriveTopicId(descriptor,\ selfRegion?)}
→
\texttt{Promise\textless{}string\textgreater{}}}{deriveTopicId(descriptor, selfRegion?) → Promise\textless string\textgreater{}}}\label{derivetopiciddescriptor-selfregion-promisestring}

Compute the 66-hex topic ID for a descriptor --- the same function the
peer uses internally. This is what you hand to someone so they can
subscribe.

{\def\LTcaptype{none} % do not increment counter
\begin{longtable}[]{@{}
  >{\raggedright\arraybackslash}p{(\linewidth - 4\tabcolsep) * \real{0.3333}}
  >{\raggedright\arraybackslash}p{(\linewidth - 4\tabcolsep) * \real{0.3333}}
  >{\raggedright\arraybackslash}p{(\linewidth - 4\tabcolsep) * \real{0.3333}}@{}}
\toprule\noalign{}
\begin{minipage}[b]{\linewidth}\raggedright
Param
\end{minipage} & \begin{minipage}[b]{\linewidth}\raggedright
Type
\end{minipage} & \begin{minipage}[b]{\linewidth}\raggedright
Notes
\end{minipage} \\
\midrule\noalign{}
\endhead
\bottomrule\noalign{}
\endlastfoot
\texttt{descriptor} & \texttt{TopicDescriptor} &
\texttt{\{\ region?,\ owner?,\ name,\ write?\ \}}. \\
\texttt{selfRegion} &
\texttt{string\ \textbackslash{}\textbar{}\ number} & fallback region
when \texttt{descriptor.region} is omitted (the publisher's node
region). \\
\end{longtable}
}

\begin{Shaded}
\begin{Highlighting}[]
\ImportTok{import}\NormalTok{ \{ deriveTopicId \} }\ImportTok{from} \StringTok{\textquotesingle{}@axona/protocol\textquotesingle{}}\OperatorTok{;}

\KeywordTok{const}\NormalTok{ lobbyId }\OperatorTok{=} \ControlFlowTok{await} \FunctionTok{deriveTopicId}\NormalTok{(\{ }\DataTypeTok{region}\OperatorTok{:} \StringTok{\textquotesingle{}useast\textquotesingle{}}\OperatorTok{,} \DataTypeTok{name}\OperatorTok{:} \StringTok{\textquotesingle{}lobby\textquotesingle{}}\NormalTok{ \})}\OperatorTok{;}
\KeywordTok{const}\NormalTok{ feedId  }\OperatorTok{=} \ControlFlowTok{await} \FunctionTok{deriveTopicId}\NormalTok{(\{ }\DataTypeTok{region}\OperatorTok{:} \StringTok{\textquotesingle{}useast\textquotesingle{}}\OperatorTok{,} \DataTypeTok{owner}\OperatorTok{:}\NormalTok{ me}\OperatorTok{.}\AttributeTok{authorId}\OperatorTok{,} \DataTypeTok{name}\OperatorTok{:} \StringTok{\textquotesingle{}profile\textquotesingle{}}\NormalTok{ \})}\OperatorTok{;}
\CommentTok{// → "89…"  (66 hex chars; leading "89" is the region byte) — share out{-}of{-}band}
\end{Highlighting}
\end{Shaded}

\textbf{Throws} \texttt{TypeError} (empty \texttt{name}),
\texttt{RangeError} (unknown region, bad \texttt{owner} format, or no
region and no \texttt{selfRegion}).

\paragraph{\texorpdfstring{\texttt{deriveTopicIdBig(descriptor,\ selfRegion?)}
→
\texttt{Promise\textless{}bigint\textgreater{}}}{deriveTopicIdBig(descriptor, selfRegion?) → Promise\textless bigint\textgreater{}}}\label{derivetopicidbigdescriptor-selfregion-promisebigint}

BigInt variant of \texttt{deriveTopicId} (the form kernel internals
hold). Same parameters and errors.

\paragraph{\texorpdfstring{\texttt{resolveTopic(descriptor,\ selfRegion?)}
→ \texttt{Promise\textless{}object\textgreater{}}
\emph{(advanced)}}{resolveTopic(descriptor, selfRegion?) → Promise\textless object\textgreater{} (advanced)}}\label{resolvetopicdescriptor-selfregion-promiseobject-advanced}

The full resolver behind \texttt{deriveTopicId}. Returns the normalized
descriptor plus the id:
\texttt{\{\ region\ (code),\ owner,\ name,\ write,\ topicId\ \}}. Use it
when you need the resolved \texttt{write}/\texttt{region}, not just the
id.

\begin{Shaded}
\begin{Highlighting}[]
\KeywordTok{const}\NormalTok{ r }\OperatorTok{=} \ControlFlowTok{await} \FunctionTok{resolveTopic}\NormalTok{(\{ }\DataTypeTok{region}\OperatorTok{:} \StringTok{\textquotesingle{}useast\textquotesingle{}}\OperatorTok{,} \DataTypeTok{owner}\OperatorTok{:}\NormalTok{ me}\OperatorTok{.}\AttributeTok{authorId}\OperatorTok{,} \DataTypeTok{name}\OperatorTok{:} \StringTok{\textquotesingle{}profile\textquotesingle{}}\NormalTok{ \})}\OperatorTok{;}
\CommentTok{// \{ region: 137, owner: \textquotesingle{}\textless{}64{-}hex\textgreater{}\textquotesingle{}, name: \textquotesingle{}profile\textquotesingle{}, write: \textquotesingle{}owner\textquotesingle{}, topicId: \textquotesingle{}89…\textquotesingle{} \}}
\end{Highlighting}
\end{Shaded}

\paragraph{\texorpdfstring{\texttt{canonical(value)} → \texttt{string}
\emph{(advanced)}}{canonical(value) → string (advanced)}}\label{canonicalvalue-string-advanced}

Stable, total, JSON-valid canonical encoding (object keys sorted at
every level). Two semantically-identical values canonicalize to the same
bytes across runs and implementations --- the basis of content-addressed
\texttt{msgId}s and signature stability.

\paragraph{\texorpdfstring{\texttt{sha256Hex(input)} →
\texttt{Promise\textless{}string\textgreater{}}
\emph{(advanced)}}{sha256Hex(input) → Promise\textless string\textgreater{} (advanced)}}\label{sha256hexinput-promisestring-advanced}

\texttt{SHA-256} of a string or \texttt{Uint8Array}, returned as
lowercase hex.

\subsubsection{4.2 publish / subscribe}\label{publish-subscribe}

\paragraph{\texorpdfstring{\texttt{peer.pub(topic,\ message,\ \{\ signWith\ \})}
→
\texttt{Promise\textless{}string\textgreater{}}}{peer.pub(topic, message, \{ signWith \}) → Promise\textless string\textgreater{}}}\label{peer.pubtopic-message-signwith-promisestring}

Publish a message to a topic. \textbf{Returns} the 64-char hex
\texttt{msgId}.

{\def\LTcaptype{none} % do not increment counter
\begin{longtable}[]{@{}
  >{\raggedright\arraybackslash}p{(\linewidth - 4\tabcolsep) * \real{0.3333}}
  >{\raggedright\arraybackslash}p{(\linewidth - 4\tabcolsep) * \real{0.3333}}
  >{\raggedright\arraybackslash}p{(\linewidth - 4\tabcolsep) * \real{0.3333}}@{}}
\toprule\noalign{}
\begin{minipage}[b]{\linewidth}\raggedright
Arg
\end{minipage} & \begin{minipage}[b]{\linewidth}\raggedright
Type
\end{minipage} & \begin{minipage}[b]{\linewidth}\raggedright
Notes
\end{minipage} \\
\midrule\noalign{}
\endhead
\bottomrule\noalign{}
\endlastfoot
\texttt{topic} & \texttt{TopicDescriptor} & \textbf{must be a
descriptor} --- a bare id is rejected (the storing node needs the
descriptor to verify the write policy). \\
\texttt{message} & \texttt{any} & JSON-serializable payload. \\
\texttt{opts.signWith} &
\texttt{AuthorIdentity\ \textbackslash{}\textbar{}\ ANONYMOUS} &
\textbf{required} --- name a signer, or pass the \texttt{ANONYMOUS}
sentinel to publish unsigned. There is no default author. \\
\end{longtable}
}

\begin{Shaded}
\begin{Highlighting}[]
\ImportTok{import}\NormalTok{ \{ createAuthorIdentity}\OperatorTok{,}\NormalTok{ ANONYMOUS \} }\ImportTok{from} \StringTok{\textquotesingle{}@axona/protocol\textquotesingle{}}\OperatorTok{;}
\KeywordTok{const}\NormalTok{ me }\OperatorTok{=} \ControlFlowTok{await} \FunctionTok{createAuthorIdentity}\NormalTok{()}\OperatorTok{;}

\KeywordTok{const}\NormalTok{ id }\OperatorTok{=} \ControlFlowTok{await}\NormalTok{ peer}\OperatorTok{.}\FunctionTok{pub}\NormalTok{(\{ }\DataTypeTok{region}\OperatorTok{:} \StringTok{\textquotesingle{}useast\textquotesingle{}}\OperatorTok{,} \DataTypeTok{name}\OperatorTok{:} \StringTok{\textquotesingle{}lobby\textquotesingle{}}\NormalTok{ \}}\OperatorTok{,}\NormalTok{ \{ }\DataTypeTok{text}\OperatorTok{:} \StringTok{\textquotesingle{}hi\textquotesingle{}}\NormalTok{ \}}\OperatorTok{,}\NormalTok{ \{ }\DataTypeTok{signWith}\OperatorTok{:}\NormalTok{ me \})}\OperatorTok{;}
\ControlFlowTok{await}\NormalTok{ peer}\OperatorTok{.}\FunctionTok{pub}\NormalTok{(\{ }\DataTypeTok{region}\OperatorTok{:} \StringTok{\textquotesingle{}useast\textquotesingle{}}\OperatorTok{,} \DataTypeTok{name}\OperatorTok{:} \StringTok{\textquotesingle{}lobby\textquotesingle{}}\NormalTok{ \}}\OperatorTok{,} \StringTok{\textquotesingle{}anon msg\textquotesingle{}}\OperatorTok{,}\NormalTok{ \{ }\DataTypeTok{signWith}\OperatorTok{:}\NormalTok{ ANONYMOUS \})}\OperatorTok{;}
\end{Highlighting}
\end{Shaded}

For an \textbf{owned} topic, \texttt{signWith} must be the owner key ---
publishing with any other key throws before it leaves the peer.

\textbf{Maximum size.} The cap is on the \emph{serialized envelope} and
defaults to the WebRTC-interoperable reliable-delivery floor (16 KiB) so
any peer on any path can receive it. Chunk larger payloads
(\texttt{@axona/protocol/std/chunk}) or publish a content-reference and
transfer bytes out-of-band.

\begin{quote}
\textbf{Chunking is reliable by default (\texttt{std/chunk}, kernel ≥
3.6.0).}
\texttt{publishChunkedBytes(peer,\ bytes,\ \{\ topic,\ signWith,\ name,\ mime\ \})}
splits a payload into floor-sized messages, publishes them, then
\textbf{verifies what the mesh actually cached and re-publishes any
gaps} --- so a \emph{reload} subscriber can reassemble from replay.
\texttt{peer.pub} is fire-and-forget into the transport buffer, so a
fast burst would otherwise drop chunks before they cache; the
verify+repair pass means you do \textbf{not} hand-tune a publish
throttle. \texttt{receiveChunkedBytes(peer,\ topic,\ \{\ timeoutMs\ \})}
reassembles and resolves with
\texttt{\{\ bytes,\ name,\ mime,\ size,\ meta\ \}}, or \textbf{rejects}
(naming the missing indices) on timeout --- it never hangs. Files are
capped to the per-topic replay-cache ceiling so you never create a
transfer a reload joiner can't complete.
\end{quote}

\begin{quote}
\textbf{Delivery (4.x).} A publish routes to the topic's K-closest
\textbf{cohort} and is replicated across it (not a single root), so it
survives the closest node churning out. A publish is still one-shot
fire-and-forget --- there is \textbf{no delivery ack to the publisher}
(the transport identity is deliberately unlinked from the author
identity) --- but two automatic mechanisms cover the gaps: a
freshly-joined (``cold'') node re-sends its first publishes a few times
over \textasciitilde1 s (v4.11.0) so a not-yet-warm routing table
doesn't strand them, and a background retry re-sends a recent publish
toward the true root until the publisher observes its own \texttt{msgId}
land. You do not tune any of this. Since v4.22.0 the cohort also
converges to the \textbf{union} of its members' history across root
transitions, so a \texttt{since:\textquotesingle{}all\textquotesingle{}}
subscriber attaching right after the root moved still replays the full
timeline, not just the newer half.
\end{quote}

\textbf{Throws} \texttt{PublishError}:

\begin{itemize}
\tightlist
\item
  \texttt{PUBLISH\_INVALID\_TOPIC} --- a bare id was passed, or the
  topic isn't a \texttt{\{\ name,\ …\ \}} object.
\item
  \texttt{TOPIC\_REGION\_REQUIRED} --- no region named and the peer has
  no node region to default to.
\item
  \texttt{PUBLISH\_NO\_PUBLISH\_IDENTITY} --- no \texttt{signWith}
  given.
\item
  \texttt{WRITE\_POLICY\_VIOLATION} --- owned topic, signer is not the
  owner.
\item
  \texttt{PUBLISH\_SIGN\_FAILED} --- \texttt{signWith} lacks
  \texttt{privateKey}/\texttt{pubkeyHex}, or signing failed.
\item
  \texttt{PUBLISH\_INVALID\_MESSAGE} --- message isn't
  JSON-serializable.
\item
  \texttt{PUBLISH\_PAYLOAD\_TOO\_LARGE} --- enveloped message exceeds
  the cap.
\end{itemize}

\paragraph{\texorpdfstring{\texttt{peer.sub(topic,\ handler,\ \{\ since?\ \})}
→
\texttt{Promise\textless{}Subscription\textgreater{}}}{peer.sub(topic, handler, \{ since? \}) → Promise\textless Subscription\textgreater{}}}\label{peer.subtopic-handler-since-promisesubscription}

Subscribe. \texttt{handler} is invoked with the full
\hyperref[envelope-type]{\texttt{Envelope}} for each delivery.

{\def\LTcaptype{none} % do not increment counter
\begin{longtable}[]{@{}
  >{\raggedright\arraybackslash}p{(\linewidth - 4\tabcolsep) * \real{0.3333}}
  >{\raggedright\arraybackslash}p{(\linewidth - 4\tabcolsep) * \real{0.3333}}
  >{\raggedright\arraybackslash}p{(\linewidth - 4\tabcolsep) * \real{0.3333}}@{}}
\toprule\noalign{}
\begin{minipage}[b]{\linewidth}\raggedright
Arg
\end{minipage} & \begin{minipage}[b]{\linewidth}\raggedright
Type
\end{minipage} & \begin{minipage}[b]{\linewidth}\raggedright
Notes
\end{minipage} \\
\midrule\noalign{}
\endhead
\bottomrule\noalign{}
\endlastfoot
\texttt{topic} &
\texttt{TopicDescriptor\ \textbackslash{}\textbar{}\ string} & a
descriptor \textbf{or} a 66-hex topic ID (the shareable read handle). \\
\texttt{handler} & \texttt{(envelope)\ =\textgreater{}\ void} & called
per delivery; also receives retraction markers
\texttt{\{\ deleted:\ true,\ msgId,\ topic\ \}}. \\
\texttt{opts.since} &
\texttt{\textquotesingle{}all\textquotesingle{}\ \textbackslash{}\textbar{}\ \textquotesingle{}latest\textquotesingle{}\ \textbackslash{}\textbar{}\ number}
& replay control (below). \\
\end{longtable}
}

\textbf{\texttt{since} modes:}

{\def\LTcaptype{none} % do not increment counter
\begin{longtable}[]{@{}
  >{\raggedright\arraybackslash}p{(\linewidth - 2\tabcolsep) * \real{0.5000}}
  >{\raggedright\arraybackslash}p{(\linewidth - 2\tabcolsep) * \real{0.5000}}@{}}
\toprule\noalign{}
\begin{minipage}[b]{\linewidth}\raggedright
\texttt{since}
\end{minipage} & \begin{minipage}[b]{\linewidth}\raggedright
What you get
\end{minipage} \\
\midrule\noalign{}
\endhead
\bottomrule\noalign{}
\endlastfoot
omitted / \texttt{undefined} & live tail only --- future messages. \\
\texttt{\textquotesingle{}latest\textquotesingle{}} & the most recent
cached message, then live tail. \\
\texttt{\textquotesingle{}all\textquotesingle{}} & everything in the
replay cache, then live tail. \\
\texttt{\textless{}number\textgreater{}} & messages newer than the
timestamp (ms epoch), then live tail. \\
\end{longtable}
}

\textbf{Returns} a \hyperref[subscription-type]{\texttt{Subscription}};
call \texttt{.stop()} to cancel.

\begin{Shaded}
\begin{Highlighting}[]
\KeywordTok{const}\NormalTok{ sub }\OperatorTok{=} \ControlFlowTok{await}\NormalTok{ peer}\OperatorTok{.}\FunctionTok{sub}\NormalTok{(\{ }\DataTypeTok{region}\OperatorTok{:} \StringTok{\textquotesingle{}useast\textquotesingle{}}\OperatorTok{,} \DataTypeTok{name}\OperatorTok{:} \StringTok{\textquotesingle{}lobby\textquotesingle{}}\NormalTok{ \}}\OperatorTok{,}\NormalTok{ (env) }\KeywordTok{=\textgreater{}}\NormalTok{ \{}
  \ControlFlowTok{if}\NormalTok{ (env}\OperatorTok{.}\AttributeTok{deleted}\NormalTok{) \{ }\FunctionTok{dropFromUI}\NormalTok{(env}\OperatorTok{.}\AttributeTok{msgId}\NormalTok{)}\OperatorTok{;} \ControlFlowTok{return}\OperatorTok{;}\NormalTok{ \}}
  \BuiltInTok{console}\OperatorTok{.}\FunctionTok{log}\NormalTok{(env}\OperatorTok{.}\AttributeTok{signerPubkey}\OperatorTok{,}\NormalTok{ env}\OperatorTok{.}\AttributeTok{message}\NormalTok{)}\OperatorTok{;}
\NormalTok{\}}\OperatorTok{,}\NormalTok{ \{ }\DataTypeTok{since}\OperatorTok{:} \StringTok{\textquotesingle{}all\textquotesingle{}}\NormalTok{ \})}\OperatorTok{;}

\CommentTok{// or subscribe by shared id:}
\KeywordTok{const}\NormalTok{ sub2 }\OperatorTok{=} \ControlFlowTok{await}\NormalTok{ peer}\OperatorTok{.}\FunctionTok{sub}\NormalTok{(lobbyId}\OperatorTok{,}\NormalTok{ onMsg}\OperatorTok{,}\NormalTok{ \{ }\DataTypeTok{since}\OperatorTok{:} \StringTok{\textquotesingle{}latest\textquotesingle{}}\NormalTok{ \})}\OperatorTok{;}

\ControlFlowTok{await}\NormalTok{ sub}\OperatorTok{.}\FunctionTok{stop}\NormalTok{()}\OperatorTok{;}
\end{Highlighting}
\end{Shaded}

\textbf{Throws} \texttt{SubscribeError(SUBSCRIBE\_HANDLER\_MISSING)}
(handler not a function); \texttt{PublishError(PUBLISH\_INVALID\_TOPIC)}
(a string that isn't a valid 66-hex id).

\paragraph{\texorpdfstring{\texttt{peer.unsub(topic)} →
\texttt{Promise\textless{}\{\ ok,\ removed\ \}\textgreater{}}}{peer.unsub(topic) → Promise\textless\{ ok, removed \}\textgreater{}}}\label{peer.unsubtopic-promise-ok-removed}

Unsubscribe by topic --- the counterpart to \texttt{sub}. Stops
\textbf{every} local subscription this peer holds for the topic (no need
to keep the handle) and, once the last one goes, sends the network
unsubscribe so the topic's roots drop this peer. Self-only by
construction. Idempotent --- returns
\texttt{\{\ ok:\ true,\ removed:\ 0\ \}} for a topic you're not on.

Takes a \textbf{descriptor} (it derives the topicId the same way
\texttt{sub}'s descriptor form does).

\begin{Shaded}
\begin{Highlighting}[]
\KeywordTok{const}\NormalTok{ \{ removed \} }\OperatorTok{=} \ControlFlowTok{await}\NormalTok{ peer}\OperatorTok{.}\FunctionTok{unsub}\NormalTok{(\{ }\DataTypeTok{region}\OperatorTok{:} \StringTok{\textquotesingle{}useast\textquotesingle{}}\OperatorTok{,} \DataTypeTok{name}\OperatorTok{:} \StringTok{\textquotesingle{}lobby\textquotesingle{}}\NormalTok{ \})}\OperatorTok{;}
\end{Highlighting}
\end{Shaded}

\subsubsection{4.3 pull / metrics}\label{pull-metrics}

\paragraph{\texorpdfstring{\texttt{peer.pull(msgId,\ \{\ topic,\ timeoutMs?\ \})}
→
\texttt{Promise\textless{}Envelope\ \textbar{}\ null\textgreater{}}}{peer.pull(msgId, \{ topic, timeoutMs? \}) → Promise\textless Envelope \textbar{} null\textgreater{}}}\label{peer.pullmsgid-topic-timeoutms-promiseenvelope-null}

Fetch one message from the topic's replay cache. Returns \texttt{null}
on a cache miss (including a message that aged out of the hold window)
--- that's expected, not an error.

{\def\LTcaptype{none} % do not increment counter
\begin{longtable}[]{@{}
  >{\raggedright\arraybackslash}p{(\linewidth - 4\tabcolsep) * \real{0.3333}}
  >{\raggedright\arraybackslash}p{(\linewidth - 4\tabcolsep) * \real{0.3333}}
  >{\raggedright\arraybackslash}p{(\linewidth - 4\tabcolsep) * \real{0.3333}}@{}}
\toprule\noalign{}
\begin{minipage}[b]{\linewidth}\raggedright
Arg
\end{minipage} & \begin{minipage}[b]{\linewidth}\raggedright
Type
\end{minipage} & \begin{minipage}[b]{\linewidth}\raggedright
Notes
\end{minipage} \\
\midrule\noalign{}
\endhead
\bottomrule\noalign{}
\endlastfoot
\texttt{msgId} & \texttt{string\ \textbackslash{}\textbar{}\ null} &
64-char hex, \textbf{or \texttt{null}/omitted for the topic's
most-recent message}. \\
\texttt{opts.topic} &
\texttt{TopicDescriptor\ \textbackslash{}\textbar{}\ string} &
descriptor \textbf{or} 66-hex id. \\
\texttt{opts.timeoutMs} & \texttt{number} & default \texttt{1000}. \\
\end{longtable}
}

\textbf{Nearest-replica reads (v4.11.1+).} A pull is answered by the
\textbf{first replica the request reaches} --- a cohort member, a child
relay, or a \texttt{host()} node --- not necessarily the topic's root.
This lowers latency, spreads reads off the root, and lets a pull that
would otherwise strand toward the root be served by any cache-holder it
passes. Two consistency tiers:

\begin{itemize}
\tightlist
\item
  \textbf{\texttt{pull(msgId)}} --- \emph{exact}.
  \texttt{msgId\ =\ H(publisher‖message)}, so a nearer copy \textbf{is}
  the copy; you always get that specific immutable message (or
  \texttt{null}).
\item
  \textbf{pull-latest} (\texttt{msgId} null) --- \emph{recent,
  eventually-consistent}. You get whatever newest that first replica
  holds, which may be a beat behind the root's very newest until the
  cohort converges. This is deliberate: it keeps a heavily-polled
  ``current value'' read from making the root a throughput bottleneck.
  If you need the linearizable newest, pull a specific \texttt{msgId}.
\end{itemize}

A cache-holding replica has not tombstoned the message (a \texttt{kill}
drops it from cache), so a pull never returns a killed message. A
successful pull slides the message's hold time forward (bounded by the
48 h ceiling).

\begin{Shaded}
\begin{Highlighting}[]
\KeywordTok{const}\NormalTok{ env    }\OperatorTok{=} \ControlFlowTok{await}\NormalTok{ peer}\OperatorTok{.}\FunctionTok{pull}\NormalTok{(msgId}\OperatorTok{,}\NormalTok{ \{ }\DataTypeTok{topic}\OperatorTok{:}\NormalTok{ feedId \})}\OperatorTok{;}
\KeywordTok{const}\NormalTok{ latest }\OperatorTok{=} \ControlFlowTok{await}\NormalTok{ peer}\OperatorTok{.}\FunctionTok{pull}\NormalTok{(}\KeywordTok{null}\OperatorTok{,}\NormalTok{ \{ }\DataTypeTok{topic}\OperatorTok{:}\NormalTok{ feedId \})}\OperatorTok{;}   \CommentTok{// most recent on the topic}
\end{Highlighting}
\end{Shaded}

\textbf{Throws} \texttt{PullError}:

\begin{itemize}
\tightlist
\item
  \texttt{PULL\_INVALID\_MSGID} --- a non-null \texttt{msgId} that isn't
  64-char hex.
\item
  \texttt{PULL\_AXONS\_UNREACHABLE} --- the manager can't service the
  request.
\end{itemize}

\paragraph{\texorpdfstring{\texttt{peer.metrics(topic,\ \{\ timeoutMs?\ \})}
→
\texttt{Promise\textless{}metricsObj\textgreater{}}}{peer.metrics(topic, \{ timeoutMs? \}) → Promise\textless metricsObj\textgreater{}}}\label{peer.metricstopic-timeoutms-promisemetricsobj}

A one-shot read of a topic's latest published metric snapshot --- for
\textbf{any} topic, open or owned. Metrics are \textbf{demand-driven}
(v4.12.0): subscribing to \texttt{metricTopic(T)} --- which
\texttt{metrics()} does internally --- sends a renewable
\emph{metrics-on} lease toward the data topic's root, and \textbf{any
node that roots the topic} (relay or ordinary client peer alike) then
publishes a signed snapshot to that \emph{open} metric topic every
\textasciitilde20 s while at least one metric subscriber remains.
\texttt{metrics()} collects whatever arrives (or is replayed from the
cache) during its short window and returns the freshest view. Owned
topics' metrics are \textbf{public too} --- anyone who can derive the id
can read them --- so there is no owner gate.

\begin{quote}
\textbf{Cold-topic note.} Publication starts \emph{on demand}, and since
v4.16.1 the root answers a freshly armed lease \textbf{immediately} ---
the first snapshot rides back at routing latency (\textasciitilde0.3 s
measured on testnet), normally inside the default 1500 ms collection
window. Under churn (the root re-homing at that instant) it can take a
few seconds; if the window closes empty the call returns
\texttt{stale:\ true} --- retry, widen the window, or --- better ---
keep a standing \texttt{sub(metricTopic(T),\ …)} and treat metrics as
the stream it is. If \emph{any} peer watched the topic's metrics within
the 48 h hold window, replay serves the last cached snapshot immediately
regardless.
\end{quote}

\textbf{Cohort-aware (v4.10.1).} Under the K-closest cohort model every
co-hosting root publishes its own snapshot, so \texttt{metrics()}
collects them over the window and \textbf{aggregates}:
\texttt{subscribers} is \textbf{summed} (each root reports only its own
subscriber subset, so the sum is the topic-wide total), while
\texttt{current\_count}, \texttt{seq}, and \texttt{bytes} are
\textbf{maxed} (they converge across the cohort via anti-entropy; max
tolerates a lagging member). \texttt{cohortSize} tells you how many
roots reported.

{\def\LTcaptype{none} % do not increment counter
\begin{longtable}[]{@{}
  >{\raggedright\arraybackslash}p{(\linewidth - 4\tabcolsep) * \real{0.3333}}
  >{\raggedright\arraybackslash}p{(\linewidth - 4\tabcolsep) * \real{0.3333}}
  >{\raggedright\arraybackslash}p{(\linewidth - 4\tabcolsep) * \real{0.3333}}@{}}
\toprule\noalign{}
\begin{minipage}[b]{\linewidth}\raggedright
Arg
\end{minipage} & \begin{minipage}[b]{\linewidth}\raggedright
Type
\end{minipage} & \begin{minipage}[b]{\linewidth}\raggedright
Notes
\end{minipage} \\
\midrule\noalign{}
\endhead
\bottomrule\noalign{}
\endlastfoot
\texttt{topic} &
\texttt{TopicDescriptor\ \textbackslash{}\textbar{}\ string} & the
\textbf{data} topic --- descriptor \textbf{or} 66-hex id. \\
\texttt{opts.timeoutMs} & \texttt{number} & how long to collect cohort
snapshots; default \texttt{1500}. \\
\end{longtable}
}

\textbf{Returns:}

\begin{Shaded}
\begin{Highlighting}[]
\NormalTok{\{}
\NormalTok{  current\_count}\OperatorTok{:} \DataTypeTok{number}\OperatorTok{;}       \CommentTok{// live (non{-}expired, non{-}killed) messages retained now — max across cohort}
\NormalTok{  seq}\OperatorTok{:}           \DataTypeTok{number}\OperatorTok{;}       \CommentTok{// dense message counter: monotonic high{-}water of total events ever}
                               \CommentTok{//   emitted on the topic (kills included) — max across cohort}
\NormalTok{  subscribers}\OperatorTok{:}   \DataTypeTok{number}\OperatorTok{;}       \CommentTok{// topic{-}wide subscriber total — summed across the cohort}
\NormalTok{  bytes}\OperatorTok{:}         \DataTypeTok{number}\OperatorTok{;}       \CommentTok{// live cached envelope bytes — max across cohort}
\NormalTok{  publishes}\OperatorTok{:}     \DataTypeTok{number}\OperatorTok{;}       \CommentTok{// present only if the publisher tracks it; else 0}
\NormalTok{  ts}\OperatorTok{:}            \DataTypeTok{number}\OperatorTok{|}\DataTypeTok{null}\OperatorTok{;}  \CommentTok{// freshest snapshot timestamp}
\NormalTok{  signer}\OperatorTok{:}        \DataTypeTok{string}\OperatorTok{|}\DataTypeTok{null}\OperatorTok{;}  \CommentTok{// the freshest snapshot envelope\textquotesingle{}s signerPubkey (provenance)}
\NormalTok{  cohortSize}\OperatorTok{:}    \DataTypeTok{number}\OperatorTok{;}       \CommentTok{// \# of distinct roots that reported a snapshot}
\NormalTok{  stale}\OperatorTok{:}         \DataTypeTok{boolean}\OperatorTok{;}      \CommentTok{// true ⇒ no snapshot seen (no metrics publisher roots this topic)}
\NormalTok{\}}
\end{Highlighting}
\end{Shaded}

Use \texttt{seq} for total-ever (published + killed),
\texttt{current\_count} for currently-live, and their gap to spot churn.
\texttt{subscribers} is a topic-wide count for UX, not billing. All
values are advisory.

\begin{Shaded}
\begin{Highlighting}[]
\KeywordTok{const}\NormalTok{ m }\OperatorTok{=} \ControlFlowTok{await}\NormalTok{ peer}\OperatorTok{.}\FunctionTok{metrics}\NormalTok{(\{ }\DataTypeTok{region}\OperatorTok{:} \StringTok{\textquotesingle{}useast\textquotesingle{}}\OperatorTok{,} \DataTypeTok{name}\OperatorTok{:} \StringTok{\textquotesingle{}lobby\textquotesingle{}}\NormalTok{ \})}\OperatorTok{;}
\ControlFlowTok{if}\NormalTok{ (}\OperatorTok{!}\NormalTok{m}\OperatorTok{.}\AttributeTok{stale}\NormalTok{) }\BuiltInTok{console}\OperatorTok{.}\FunctionTok{log}\NormalTok{(}\VerbatimStringTok{\textasciigrave{}}\SpecialCharTok{$\{}\NormalTok{m}\OperatorTok{.}\AttributeTok{subscribers}\SpecialCharTok{\}}\VerbatimStringTok{ subscribers, }\SpecialCharTok{$\{}\NormalTok{m}\OperatorTok{.}\AttributeTok{current\_count}\SpecialCharTok{\}}\VerbatimStringTok{ live (}\SpecialCharTok{$\{}\NormalTok{m}\OperatorTok{.}\AttributeTok{seq}\SpecialCharTok{\}}\VerbatimStringTok{ total), }\SpecialCharTok{$\{}\NormalTok{m}\OperatorTok{.}\AttributeTok{cohortSize}\SpecialCharTok{\}}\VerbatimStringTok{ roots\textasciigrave{}}\NormalTok{)}\OperatorTok{;}
\end{Highlighting}
\end{Shaded}

\begin{quote}
\textbf{For a live dashboard, prefer \texttt{sub(metricTopic(T),\ …)}
directly} (below) --- one standing subscription gives you the latest
snapshot plus a rolling history. \texttt{metrics()} is the convenience
one-shot. Trust is \textbf{advisory}: the metric topic is open, so check
\texttt{signer} if you need provenance (pin to a known relay key).
\end{quote}

\paragraph{\texorpdfstring{\texttt{metricTopic(dataTopicId)} →
\texttt{TopicDescriptor} \emph{(subscribe to metrics --- the live
path)}}{metricTopic(dataTopicId) → TopicDescriptor (subscribe to metrics --- the live path)}}\label{metrictopicdatatopicid-topicdescriptor-subscribe-to-metrics-the-live-path}

The producer side of the convention. Compute the derived, open metric
topic with \texttt{metricTopic()} and \texttt{sub()} it --- \textbf{your
subscription is what turns publishing on}: it routes a renewable
\emph{metrics-on} lease to the data topic's root, and while the lease is
fresh the root publishes a signed snapshot every \textasciitilde20 s,
for \textbf{both open and owned} data topics. You get the latest
snapshot via replay-on-subscribe plus every update --- one subscription
instead of a poll --- and, because snapshots are ordinary messages that
age out at the 48 h hold ceiling, a \textbf{rolling \textasciitilde48 h
history for free} to plot trends. This is also how you subscribe to an
\textbf{owned} topic's metrics without owning it.

\textbf{When snapshots arrive.} The timing contract (since v4.16.1):

\begin{itemize}
\tightlist
\item
  \textbf{First snapshot: at routing latency.} The root answers the
  moment your subscription's lease arms --- \textbf{\textasciitilde0.3 s
  measured on testnet}, arriving with (or before) your data-topic
  replay, whether or not anyone has ever published or watched. Allow a
  few seconds if the root churns at exactly that moment (it must re-home
  first). If any peer watched this topic's metrics within the last 48 h,
  replay hands you the most recent snapshot immediately as well.
\item
  \textbf{Cadence:} one snapshot every \textbf{\textasciitilde20 s} per
  rooting node, for as long as at least one metric subscriber remains.
\item
  \textbf{Shut-off:} the lease lapses \textbf{\textasciitilde70 s} after
  the last metric subscriber unsubscribes; publishing stops until demand
  returns.
\item
  \textbf{Silence is \emph{unknown}, not zero.} Until the first snapshot
  (or your data replay) arrives, don't render ``no activity'' as a
  definitive answer; a \texttt{current\_count:\ 0} snapshot is the real
  ``nothing here.''
\end{itemize}

\begin{Shaded}
\begin{Highlighting}[]
\ImportTok{import}\NormalTok{ \{ deriveTopicId}\OperatorTok{,}\NormalTok{ metricTopic \} }\ImportTok{from} \StringTok{\textquotesingle{}@axona/protocol\textquotesingle{}}\OperatorTok{;}

\KeywordTok{const}\NormalTok{ id }\OperatorTok{=} \ControlFlowTok{await} \FunctionTok{deriveTopicId}\NormalTok{(\{ }\DataTypeTok{region}\OperatorTok{:} \StringTok{\textquotesingle{}useast\textquotesingle{}}\OperatorTok{,} \DataTypeTok{name}\OperatorTok{:} \StringTok{\textquotesingle{}lobby\textquotesingle{}}\NormalTok{ \})}\OperatorTok{;}
\ControlFlowTok{await}\NormalTok{ peer}\OperatorTok{.}\FunctionTok{sub}\NormalTok{(}\FunctionTok{metricTopic}\NormalTok{(id)}\OperatorTok{,}\NormalTok{ (env) }\KeywordTok{=\textgreater{}}\NormalTok{ \{}
  \KeywordTok{const}\NormalTok{ m }\OperatorTok{=} \BuiltInTok{JSON}\OperatorTok{.}\FunctionTok{parse}\NormalTok{(env}\OperatorTok{.}\AttributeTok{message}\NormalTok{)}\OperatorTok{;}
  \CommentTok{// \{ topic, ts, by, signer, current\_count, seq, subscribers, bytes \}}
  \FunctionTok{render}\NormalTok{(m}\OperatorTok{.}\AttributeTok{subscribers}\OperatorTok{,}\NormalTok{ m}\OperatorTok{.}\AttributeTok{current\_count}\NormalTok{)}\OperatorTok{;}
\NormalTok{\}}\OperatorTok{,}\NormalTok{ \{ }\DataTypeTok{since}\OperatorTok{:} \StringTok{\textquotesingle{}all\textquotesingle{}}\NormalTok{ \})}\OperatorTok{;}   \CommentTok{// since:\textquotesingle{}all\textquotesingle{} → latest snapshot + the rolling history}
\end{Highlighting}
\end{Shaded}

{\def\LTcaptype{none} % do not increment counter
\begin{longtable}[]{@{}
  >{\raggedright\arraybackslash}p{(\linewidth - 4\tabcolsep) * \real{0.3333}}
  >{\raggedright\arraybackslash}p{(\linewidth - 4\tabcolsep) * \real{0.3333}}
  >{\raggedright\arraybackslash}p{(\linewidth - 4\tabcolsep) * \real{0.3333}}@{}}
\toprule\noalign{}
\begin{minipage}[b]{\linewidth}\raggedright
name
\end{minipage} & \begin{minipage}[b]{\linewidth}\raggedright
returns
\end{minipage} & \begin{minipage}[b]{\linewidth}\raggedright
notes
\end{minipage} \\
\midrule\noalign{}
\endhead
\bottomrule\noalign{}
\endlastfoot
\texttt{metricTopic(dataTopicId)} & \texttt{TopicDescriptor} & open
topic
\texttt{\{\ region,\ name:\ \textquotesingle{}axona:metric:\textquotesingle{}\ +\ id\ \}}.
Pass the \textbf{resolved 66-hex} id
(\texttt{await\ deriveTopicId(desc)}), not the descriptor. Region byte
is inherited from the data topic. \\
\texttt{isMetricTopic(descriptor)} & \texttt{boolean} & true if a topic
is itself a metric topic (the recursion guard a relay uses). \\
\texttt{isMetricTopicName(name)} / \texttt{dataTopicIdOf(descriptor)} &
\texttt{boolean} / \texttt{string\textbackslash{}\textbar{}null} &
name-prefix test; inverse (metric topic → data id). \\
\texttt{METRIC\_NAMESPACE} & \texttt{string} & the frozen reserved
prefix \texttt{\textquotesingle{}axona:metric:\textquotesingle{}}. \\
\end{longtable}
}

Trust is \textbf{advisory}: the metric topic is open (anyone may publish
to it), so treat a snapshot as a hint --- if you need to, pin trust to a
known relay's \texttt{env.signerPubkey}. The protocol does not prove a
snapshot is authoritative.

\subsubsection{4.4 owner + creator ops:
kill}\label{owner-creator-ops-kill}

These take a \textbf{descriptor} (not a bare id) and a signer via
\texttt{\{\ signWith\ \}}.

\paragraph{\texorpdfstring{\texttt{peer.kill(topic,\ msgId,\ \{\ signWith\ \})}
→
\texttt{Promise\textless{}\{\ ok\ \}\textgreater{}}}{peer.kill(topic, msgId, \{ signWith \}) → Promise\textless\{ ok \}\textgreater{}}}\label{peer.killtopic-msgid-signwith-promise-ok}

Retract a message you published --- ``unsend''. Authorized by
\textbf{authorship}: the roots accept the kill only if its signer
matches the signer of the cached message, so you can only kill messages
\textbf{you} signed. The kill is distributed to the topic's
\textbf{K-closest cohort} (v4.10.0), not one root: every cohort member
drops it from its replay cache, tombstones the \texttt{msgId}, and every
internal history transfer (hand-off, catch-up, backup replication)
carries that tombstone --- so a late subscriber attaching to any cohort
member can't be served the killed copy. Subscribers get a delete marker
(\texttt{\{\ deleted:\ true,\ msgId,\ topic\ \}}).

{\def\LTcaptype{none} % do not increment counter
\begin{longtable}[]{@{}
  >{\raggedright\arraybackslash}p{(\linewidth - 4\tabcolsep) * \real{0.3333}}
  >{\raggedright\arraybackslash}p{(\linewidth - 4\tabcolsep) * \real{0.3333}}
  >{\raggedright\arraybackslash}p{(\linewidth - 4\tabcolsep) * \real{0.3333}}@{}}
\toprule\noalign{}
\begin{minipage}[b]{\linewidth}\raggedright
Arg
\end{minipage} & \begin{minipage}[b]{\linewidth}\raggedright
Type
\end{minipage} & \begin{minipage}[b]{\linewidth}\raggedright
Notes
\end{minipage} \\
\midrule\noalign{}
\endhead
\bottomrule\noalign{}
\endlastfoot
\texttt{topic} & \texttt{TopicDescriptor} & the topic the message was
published to. \\
\texttt{msgId} & \texttt{string} & \textbf{required} 64-char hex (the
value \texttt{pub} returned). \\
\texttt{opts.signWith} & \texttt{AuthorIdentity} & the \textbf{same
author key} that published the message. \\
\end{longtable}
}

\begin{Shaded}
\begin{Highlighting}[]
\KeywordTok{const}\NormalTok{ id }\OperatorTok{=} \ControlFlowTok{await}\NormalTok{ peer}\OperatorTok{.}\FunctionTok{pub}\NormalTok{(\{ }\DataTypeTok{region}\OperatorTok{:} \StringTok{\textquotesingle{}useast\textquotesingle{}}\OperatorTok{,} \DataTypeTok{name}\OperatorTok{:} \StringTok{\textquotesingle{}lobby\textquotesingle{}}\NormalTok{ \}}\OperatorTok{,}\NormalTok{ \{ }\DataTypeTok{text}\OperatorTok{:} \StringTok{\textquotesingle{}oops\textquotesingle{}}\NormalTok{ \}}\OperatorTok{,}\NormalTok{ \{ }\DataTypeTok{signWith}\OperatorTok{:}\NormalTok{ me \})}\OperatorTok{;}
\ControlFlowTok{await}\NormalTok{ peer}\OperatorTok{.}\FunctionTok{kill}\NormalTok{(\{ }\DataTypeTok{region}\OperatorTok{:} \StringTok{\textquotesingle{}useast\textquotesingle{}}\OperatorTok{,} \DataTypeTok{name}\OperatorTok{:} \StringTok{\textquotesingle{}lobby\textquotesingle{}}\NormalTok{ \}}\OperatorTok{,}\NormalTok{ id}\OperatorTok{,}\NormalTok{ \{ }\DataTypeTok{signWith}\OperatorTok{:}\NormalTok{ me \})}\OperatorTok{;}
\end{Highlighting}
\end{Shaded}

\begin{quote}
Best-effort redaction, not a cryptographic un-send: a subscriber who
already has the plaintext can keep it; an anonymous message has no
provable creator and \textbf{cannot} be killed.
\end{quote}

\textbf{Throws} \texttt{KillError}: \texttt{KILL\_INVALID\_MSGID} (bad
msgId), \texttt{KILL\_SIGN\_FAILED} (no/invalid \texttt{signWith}, or
signing failed). A network that can't authorize the kill is \textbf{not}
an error --- it just won't take effect.

\begin{quote}
\textbf{Removed in v4.3.0: \texttt{peer.unpub()} and
\texttt{peer.touch()}.} \texttt{kill(topic,\ msgId)} is now the single
retraction primitive --- retract individual messages you signed.
(\texttt{unpub} cleared a whole owned-topic queue; \texttt{touch} was a
hold-time keep-alive. Both are gone --- a topic's queue empties
naturally as messages age out at the 48 h hold ceiling, and a
still-relevant message is kept current by re-publishing.)
\texttt{touch()} remains callable as a no-op for source compatibility
but does nothing; do not use it.
\end{quote}

\subsubsection{4.5 hosting: host / unhost}\label{hosting-host-unhost}

\paragraph{\texorpdfstring{\texttt{peer.host(topic?,\ opts?)} →
\texttt{Promise\textless{}\{\ ok,\ scope,\ topicId?\ \}\textgreater{}}}{peer.host(topic?, opts?) → Promise\textless\{ ok, scope, topicId? \}\textgreater{}}}\label{peer.hosttopic-opts-promise-ok-scope-topicid}

Store and serve a topic for other peers \textbf{without subscribing} ---
the relay / infrastructure primitive. The node becomes a willing
root/replica (publishes land on it, subscribers pull replays from it)
but registers \textbf{no handler} and delivers nothing to a local app.

\begin{itemize}
\tightlist
\item
  \texttt{peer.host()} --- host this node's \textbf{own keyspace
  neighborhood}: get recruited as a root for whatever topics land near
  this node's id (``host whatever lands near me''). The zero-config
  relay mode; returns
  \texttt{\{\ ok:\ true,\ scope:\ \textquotesingle{}keyspace\textquotesingle{}\ \}}.
\item
  \texttt{peer.host(topic)} --- host one specific topic (descriptor).
  Returns
  \texttt{\{\ ok:\ true,\ scope:\ \textquotesingle{}topic\textquotesingle{},\ topicId\ \}}.
\end{itemize}

Wire-compatible with every kernel (reuses \texttt{subscribe-k}),
respects the proximity gate, and is idempotent.

\begin{Shaded}
\begin{Highlighting}[]
\ControlFlowTok{await}\NormalTok{ peer}\OperatorTok{.}\FunctionTok{host}\NormalTok{()}\OperatorTok{;}                                          \CommentTok{// headless relay}
\ControlFlowTok{await}\NormalTok{ peer}\OperatorTok{.}\FunctionTok{host}\NormalTok{(\{ }\DataTypeTok{region}\OperatorTok{:} \StringTok{\textquotesingle{}useast\textquotesingle{}}\OperatorTok{,} \DataTypeTok{name}\OperatorTok{:} \StringTok{\textquotesingle{}pow{-}results\textquotesingle{}}\NormalTok{ \})}\OperatorTok{;} \CommentTok{// host one topic}
\end{Highlighting}
\end{Shaded}

\paragraph{\texorpdfstring{\texttt{peer.unhost(topic?,\ opts?)} →
\texttt{Promise\textless{}\{\ ok,\ scope\ \}\textgreater{}}}{peer.unhost(topic?, opts?) → Promise\textless\{ ok, scope \}\textgreater{}}}\label{peer.unhosttopic-opts-promise-ok-scope}

Counterpart to \texttt{host}. \texttt{unhost()} turns off keyspace
hosting; \texttt{unhost(topic)} drops one hosted topic. Does
\textbf{not} touch your subscriptions. Idempotent.

\paragraph{\texorpdfstring{\texttt{peer.rootedTopics()} →
\texttt{Array\textless{}\{\ topicId,\ descriptor,\ current\_count,\ seq,\ subscribers,\ bytes\ \}\textgreater{}}
\emph{(infra)}}{peer.rootedTopics() → Array\textless\{ topicId, descriptor, current\_count, seq, subscribers, bytes \}\textgreater{} (infra)}}\label{peer.rootedtopics-array-topicid-descriptor-current_count-seq-subscribers-bytes-infra}

Synchronous, local-only introspection of the topics this node currently
\textbf{roots} --- each with its signed topic descriptor (recovered from
a cached envelope, or \texttt{null} for an empty/cold role) and a
locally-computed snapshot (\texttt{current\_count}, \texttt{seq}, this
member's \texttt{subscribers} subset, \texttt{bytes}). No network
(unlike \texttt{metrics()}). This is the read side that powers the relay
metric-publish loop (§4.3): walk it, skip \texttt{isMetricTopic(d)} and
non-open descriptors, and \texttt{pub(metricTopic(topicId),\ …)} for the
rest. Returns \texttt{{[}{]}} on a routing-only peer with no
\texttt{AxonaManager}.

\subsubsection{4.6 topic limits}\label{topic-limits}

Every topic's replay queue is bounded:

\begin{itemize}
\tightlist
\item
  \textbf{Max messages} --- default 100, max 256. When full, the
  lowest-ordered message (by signed \texttt{seq}) is evicted
  deterministically. A max of 1 is a \emph{retained/latest-value} slot;
  pair it with \texttt{peer.pull(null,\ …)}.
\item
  \textbf{Hold time} --- default 24 h, hard ceiling 48 h. A message past
  its hold is swept (stops being delivered or pulled). A \texttt{pull}
  or \texttt{touch} slides the hold forward, bounded by the ceiling.
\item
  \textbf{Per-publisher quota (open topics only)} --- one publisher
  occupies at most ¼ of the queue, so a flooder can't evict everyone
  else. Owner-gated topics aren't quota'd.
\end{itemize}

\begin{center}\rule{0.5\linewidth}{0.5pt}\end{center}

\subsection{5. Direct messaging}\label{direct-messaging}

1:1 messages between specific peers, bypassing pub/sub.
\texttt{targetId} is the 66-char hex nodeId; the peer must already be in
the synaptome.

\subsubsection{\texorpdfstring{\texttt{peer.send(targetId,\ message)} →
\texttt{Promise\textless{}reply\textgreater{}}}{peer.send(targetId, message) → Promise\textless reply\textgreater{}}}\label{peer.sendtargetid-message-promisereply}

Request/response. Resolves to whatever the recipient's
\texttt{onMessage} handler returns.

\begin{Shaded}
\begin{Highlighting}[]
\KeywordTok{const}\NormalTok{ reply }\OperatorTok{=} \ControlFlowTok{await}\NormalTok{ alice}\OperatorTok{.}\FunctionTok{send}\NormalTok{(bobIdHex}\OperatorTok{,}\NormalTok{ \{ }\DataTypeTok{kind}\OperatorTok{:} \StringTok{\textquotesingle{}ping\textquotesingle{}}\OperatorTok{,} \DataTypeTok{body}\OperatorTok{:} \StringTok{\textquotesingle{}hi\textquotesingle{}}\NormalTok{ \})}\OperatorTok{;}
\end{Highlighting}
\end{Shaded}

\textbf{Throws} \texttt{TypeError} if \texttt{targetId} isn't 66-char
hex.

\subsubsection{\texorpdfstring{\texttt{peer.notify(targetId,\ message)}
→
\texttt{Promise\textless{}void\textgreater{}}}{peer.notify(targetId, message) → Promise\textless void\textgreater{}}}\label{peer.notifytargetid-message-promisevoid}

Fire-and-forget. Resolves once enqueued, not when delivered.

\begin{Shaded}
\begin{Highlighting}[]
\ControlFlowTok{await}\NormalTok{ alice}\OperatorTok{.}\FunctionTok{notify}\NormalTok{(bobIdHex}\OperatorTok{,}\NormalTok{ \{ }\DataTypeTok{kind}\OperatorTok{:} \StringTok{\textquotesingle{}typing\textquotesingle{}}\OperatorTok{,} \DataTypeTok{user}\OperatorTok{:} \StringTok{\textquotesingle{}alice\textquotesingle{}}\NormalTok{ \})}\OperatorTok{;}
\end{Highlighting}
\end{Shaded}

\textbf{Throws} \texttt{TypeError} if \texttt{targetId} isn't 66-char
hex.

\subsubsection{\texorpdfstring{\texttt{peer.onMessage(handler)} →
\texttt{void}}{peer.onMessage(handler) → void}}\label{peer.onmessagehandler-void}

Register the single inbound direct-message handler. Calling it again
\textbf{replaces} the previous handler. Signature
\texttt{(senderId,\ message)\ =\textgreater{}\ reply\ \textbar{}\ void}:
a returned value resolves the sender's \texttt{send()}; for
\texttt{notify()} callers the return is discarded.

\begin{Shaded}
\begin{Highlighting}[]
\NormalTok{peer}\OperatorTok{.}\FunctionTok{onMessage}\NormalTok{(}\KeywordTok{async}\NormalTok{ (senderIdHex}\OperatorTok{,}\NormalTok{ message) }\KeywordTok{=\textgreater{}}\NormalTok{ \{}
  \ControlFlowTok{switch}\NormalTok{ (message}\OperatorTok{?.}\AttributeTok{kind}\NormalTok{) \{}
    \ControlFlowTok{case} \StringTok{\textquotesingle{}ping\textquotesingle{}}\OperatorTok{:}   \ControlFlowTok{return}\NormalTok{ \{ }\DataTypeTok{reply}\OperatorTok{:} \StringTok{\textquotesingle{}pong\textquotesingle{}}\NormalTok{ \}}\OperatorTok{;}   \CommentTok{// resolves send()}
    \ControlFlowTok{case} \StringTok{\textquotesingle{}typing\textquotesingle{}}\OperatorTok{:} \BuiltInTok{console}\OperatorTok{.}\FunctionTok{log}\NormalTok{(senderIdHex}\OperatorTok{,} \StringTok{\textquotesingle{}typing\textquotesingle{}}\NormalTok{)}\OperatorTok{;} \ControlFlowTok{return}\OperatorTok{;}  \CommentTok{// notify discards}
\NormalTok{  \}}
\NormalTok{\})}\OperatorTok{;}
\end{Highlighting}
\end{Shaded}

\texttt{senderId} is the transport-authenticated sender, surfaced as
66-hex.

\begin{center}\rule{0.5\linewidth}{0.5pt}\end{center}

\subsection{6. Mesh introspection +
events}\label{mesh-introspection-events}

\subsubsection{\texorpdfstring{\texttt{peer.peers()} →
\texttt{string{[}{]}}}{peer.peers() → string{[}{]}}}\label{peer.peers-string}

Current synaptome membership as 66-char hex strings.

\subsubsection{\texorpdfstring{\texttt{peer.onPeerJoin(handler)} →
\texttt{()\ =\textgreater{}\ void}}{peer.onPeerJoin(handler) → () =\textgreater{} void}}\label{peer.onpeerjoinhandler-void}

Fires when a peer is admitted to the synaptome. Handler receives
\texttt{(peerIdHex,\ event)}. Returns an unsubscribe function.

\begin{Shaded}
\begin{Highlighting}[]
\KeywordTok{const}\NormalTok{ off }\OperatorTok{=}\NormalTok{ peer}\OperatorTok{.}\FunctionTok{onPeerJoin}\NormalTok{((id}\OperatorTok{,}\NormalTok{ ev) }\KeywordTok{=\textgreater{}} \BuiltInTok{console}\OperatorTok{.}\FunctionTok{log}\NormalTok{(}\StringTok{\textquotesingle{}admitted\textquotesingle{}}\OperatorTok{,}\NormalTok{ id}\OperatorTok{,}\NormalTok{ ev}\OperatorTok{?.}\AttributeTok{addedBy}\NormalTok{))}\OperatorTok{;}
\end{Highlighting}
\end{Shaded}

\subsubsection{\texorpdfstring{\texttt{peer.onPeerLeave(handler)} →
\texttt{()\ =\textgreater{}\ void}}{peer.onPeerLeave(handler) → () =\textgreater{} void}}\label{peer.onpeerleavehandler-void}

Symmetric --- fires on eviction (TTL, churn, explicit unbind). Same
signature.

\subsubsection{\texorpdfstring{\texttt{peer.lookup(targetKey)} →
\texttt{Promise\textless{}lookupResult\textgreater{}}}{peer.lookup(targetKey) → Promise\textless lookupResult\textgreater{}}}\label{peer.lookuptargetkey-promiselookupresult}

Iterative XOR-routed walk to find \texttt{targetKey} (typically a peer
ID, but any 264-bit key works). \texttt{targetKey} is a BigInt.

\begin{Shaded}
\begin{Highlighting}[]
\KeywordTok{const}\NormalTok{ r }\OperatorTok{=} \ControlFlowTok{await}\NormalTok{ peer}\OperatorTok{.}\FunctionTok{lookup}\NormalTok{(}\BuiltInTok{BigInt}\NormalTok{(}\StringTok{\textquotesingle{}0x\textquotesingle{}} \OperatorTok{+}\NormalTok{ targetHex))}\OperatorTok{;}
\CommentTok{// \{ found: boolean, hops: number, time: number /* ms */, path: bigint[] \}}
\end{Highlighting}
\end{Shaded}

Returns \texttt{null} if the node isn't alive.

\subsubsection{\texorpdfstring{\texttt{peer.health()} →
\texttt{healthObj}}{peer.health() → healthObj}}\label{peer.health-healthobj}

A cheap synchronous diagnostic snapshot --- safe to poll on a UI tick:

\begin{Shaded}
\begin{Highlighting}[]
\NormalTok{peer}\OperatorTok{.}\FunctionTok{health}\NormalTok{()}\OperatorTok{;}
\CommentTok{// \{}
\CommentTok{//   nodeId, synaptomeSize, peers: [...], subscriptions,}
\CommentTok{//   axonRoles: [\{ topic, isRoot, children, cacheSize \}],}
\CommentTok{//   hosting:   \{ keyspace, topics \} | null,}
\CommentTok{//   wireVersion, started,}
\CommentTok{//   transport: \{ boundCount, meshChannels, meshOpen, meshBound, bridgeState \} | null,}
\CommentTok{//   meshDegraded: boolean,}
\CommentTok{// \}}
\end{Highlighting}
\end{Shaded}

\texttt{meshDegraded\ ===\ true} (open channels materially exceed
authenticated binds) is the routing-truth signal --- a single tick can
be a mid-handshake transient; a value that stays \texttt{true} across
polls is the real warning. \texttt{transport} is \texttt{null} on
sim/node transports, populated on web.

\subsubsection{\texorpdfstring{\texttt{peer.onLog(level,\ handler)} →
\texttt{()\ =\textgreater{}\ void}}{peer.onLog(level, handler) → () =\textgreater{} void}}\label{peer.onloglevel-handler-void}

Subscribe to structured log events. \textbf{The level comes first}
(\texttt{\textquotesingle{}debug\textquotesingle{}\ \textbar{}\ \textquotesingle{}info\textquotesingle{}\ \textbar{}\ \textquotesingle{}warn\textquotesingle{}\ \textbar{}\ \textquotesingle{}error\textquotesingle{}});
handler is \texttt{(msg,\ context?)\ =\textgreater{}\ void}. Returns an
unsubscribe function.

\begin{Shaded}
\begin{Highlighting}[]
\KeywordTok{const}\NormalTok{ off }\OperatorTok{=}\NormalTok{ peer}\OperatorTok{.}\FunctionTok{onLog}\NormalTok{(}\StringTok{\textquotesingle{}warn\textquotesingle{}}\OperatorTok{,}\NormalTok{ (msg}\OperatorTok{,}\NormalTok{ ctx) }\KeywordTok{=\textgreater{}} \BuiltInTok{console}\OperatorTok{.}\FunctionTok{warn}\NormalTok{(msg}\OperatorTok{,}\NormalTok{ ctx))}\OperatorTok{;}
\end{Highlighting}
\end{Shaded}

\textbf{Throws} \texttt{TypeError} on an invalid level or non-function
handler.

\subsubsection{\texorpdfstring{\texttt{peer.onError(handler)} →
\texttt{()\ =\textgreater{}\ void}}{peer.onError(handler) → () =\textgreater{} void}}\label{peer.onerrorhandler-void}

Fires on background \texttt{AxonaError} emissions (transport failures
during heartbeat, persistence warnings) the kernel surfaces
asynchronously rather than throwing. Returns an unsubscribe function.

\subsubsection{\texorpdfstring{\texttt{peer.onUpgradeRequired(handler)}
→
\texttt{()\ =\textgreater{}\ void}}{peer.onUpgradeRequired(handler) → () =\textgreater{} void}}\label{peer.onupgraderequiredhandler-void}

Fires on a wire-version handshake mismatch, with the
\texttt{UpgradeRequiredError} carrying
\texttt{\{\ reason,\ serverVersion,\ clientVersion,\ downloadUrl\ \}}.

\begin{Shaded}
\begin{Highlighting}[]
\NormalTok{peer}\OperatorTok{.}\FunctionTok{onUpgradeRequired}\NormalTok{((err) }\KeywordTok{=\textgreater{}} \FunctionTok{showUpgradeBanner}\NormalTok{(err}\OperatorTok{.}\AttributeTok{context}\OperatorTok{.}\AttributeTok{downloadUrl}\NormalTok{))}\OperatorTok{;}
\end{Highlighting}
\end{Shaded}

\begin{center}\rule{0.5\linewidth}{0.5pt}\end{center}

\subsection{7. Envelopes + verification}\label{envelopes-verification}

The kernel builds and verifies envelopes for you inside
\texttt{peer.pub} / \texttt{peer.sub}. These low-level helpers are for
verifying what you consume, signing a DM payload by hand, or testing.

\subsubsection{\texorpdfstring{\texttt{buildEnvelope(\{\ topic,\ message,\ ts?,\ seq?,\ identity,\ sign?\ \})}
→
\texttt{Promise\textless{}Envelope\textgreater{}}}{buildEnvelope(\{ topic, message, ts?, seq?, identity, sign? \}) → Promise\textless Envelope\textgreater{}}}\label{buildenvelope-topic-message-ts-seq-identity-sign-promiseenvelope}

{\def\LTcaptype{none} % do not increment counter
\begin{longtable}[]{@{}
  >{\raggedright\arraybackslash}p{(\linewidth - 4\tabcolsep) * \real{0.3333}}
  >{\raggedright\arraybackslash}p{(\linewidth - 4\tabcolsep) * \real{0.3333}}
  >{\raggedright\arraybackslash}p{(\linewidth - 4\tabcolsep) * \real{0.3333}}@{}}
\toprule\noalign{}
\begin{minipage}[b]{\linewidth}\raggedright
Param
\end{minipage} & \begin{minipage}[b]{\linewidth}\raggedright
Type
\end{minipage} & \begin{minipage}[b]{\linewidth}\raggedright
Notes
\end{minipage} \\
\midrule\noalign{}
\endhead
\bottomrule\noalign{}
\endlastfoot
\texttt{topic} & \texttt{TopicDescriptor} & the descriptor --- signed,
so the signature binds the exact write policy. \\
\texttt{message} & \texttt{any} & JSON-serializable. \\
\texttt{ts} & \texttt{number} & ms timestamp; defaults to
\texttt{Date.now()}. \\
\texttt{seq} & \texttt{number} & per-publisher monotonic sequence
(default \texttt{0}; \texttt{peer.pub} supplies a real value). \\
\texttt{identity} & \texttt{AuthorIdentity} & required when
\texttt{sign:\ true}. \\
\texttt{sign} & \texttt{boolean} & default \texttt{true}. \\
\end{longtable}
}

\textbf{Returns} an \hyperref[envelope-type]{\texttt{Envelope}}. The
signature covers a domain-tagged core
(\texttt{axona:pubsub-envelope:v2}).

\textbf{Throws} \texttt{TypeError} if \texttt{identity} lacks
\texttt{privateKey}/\texttt{pubkeyHex} when \texttt{sign:\ true}.

\subsubsection{\texorpdfstring{\texttt{verifyEnvelope(envelope)} →
\texttt{Promise\textless{}\{\ ok,\ reason?,\ signed\ \}\textgreater{}}}{verifyEnvelope(envelope) → Promise\textless\{ ok, reason?, signed \}\textgreater{}}}\label{verifyenvelopeenvelope-promise-ok-reason-signed}

Verify the signature against the domain-tagged
\texttt{(seq,\ ts,\ topic,\ message)} core and that \texttt{msgId}
matches. \textbf{Returns a result object, not a boolean} --- check
\texttt{.ok}.

\begin{Shaded}
\begin{Highlighting}[]
\KeywordTok{const}\NormalTok{ res }\OperatorTok{=} \ControlFlowTok{await} \FunctionTok{verifyEnvelope}\NormalTok{(env)}\OperatorTok{;}
\ControlFlowTok{if}\NormalTok{ (}\OperatorTok{!}\NormalTok{res}\OperatorTok{.}\AttributeTok{ok}\NormalTok{) }\BuiltInTok{console}\OperatorTok{.}\FunctionTok{warn}\NormalTok{(}\StringTok{\textquotesingle{}forged/invalid\textquotesingle{}}\OperatorTok{,}\NormalTok{ env}\OperatorTok{.}\AttributeTok{msgId}\OperatorTok{,}\NormalTok{ res}\OperatorTok{.}\AttributeTok{reason}\NormalTok{)}\OperatorTok{;}
\CommentTok{// res.signed === false for a valid unsigned envelope (res.ok === true)}
\end{Highlighting}
\end{Shaded}

\begin{quote}
A common mistake: \texttt{if\ (!(await\ verifyEnvelope(env)))} is always
false --- the return value is a truthy object. Always test \texttt{.ok}.
\end{quote}

\texttt{reason} (when \texttt{!ok}) is one of \texttt{not\_an\_object},
\texttt{missing\_msgId}, \texttt{missing\_ts}, \texttt{missing\_topic},
\texttt{missing\_message}, \texttt{missing\_seq},
\texttt{missing\_signerPubkey}, \texttt{unknown\_signature\_scheme},
\texttt{wrong\_key\_or\_signature\_length}, \texttt{bad\_signature},
\texttt{bad\_msgid}. \texttt{verifyEnvelope} checks signature + msgId
but \textbf{not} freshness (the legitimate replay path serves cached
history).

\subsubsection{\texorpdfstring{\texttt{computeMsgId(\{\ publisher?,\ message\ \})}
→
\texttt{Promise\textless{}string\textgreater{}}}{computeMsgId(\{ publisher?, message \}) → Promise\textless string\textgreater{}}}\label{computemsgid-publisher-message-promisestring}

Stand-alone msgId computation matching \texttt{buildEnvelope}:
\texttt{sha256(canonical(\{\ publisher,\ message\ \}))}, where
\texttt{publisher} is the signer's pubkey (or \texttt{null} for
anonymous). Time and seq are deliberately not folded in.

\begin{Shaded}
\begin{Highlighting}[]
\KeywordTok{const}\NormalTok{ msgId }\OperatorTok{=} \ControlFlowTok{await} \FunctionTok{computeMsgId}\NormalTok{(\{ }\DataTypeTok{publisher}\OperatorTok{:}\NormalTok{ env}\OperatorTok{.}\AttributeTok{signerPubkey}\OperatorTok{,} \DataTypeTok{message}\OperatorTok{:}\NormalTok{ env}\OperatorTok{.}\AttributeTok{message}\NormalTok{ \})}\OperatorTok{;}
\BuiltInTok{console}\OperatorTok{.}\FunctionTok{log}\NormalTok{(msgId }\OperatorTok{===}\NormalTok{ env}\OperatorTok{.}\AttributeTok{msgId}\NormalTok{)}\OperatorTok{;}   \CommentTok{// true}
\end{Highlighting}
\end{Shaded}

\subsubsection{\texorpdfstring{\texttt{checkFreshness(envelope,\ \{\ now?,\ maxSkewMs?\ \})}
→
\texttt{\{\ ok,\ reason?\ \}}}{checkFreshness(envelope, \{ now?, maxSkewMs? \}) → \{ ok, reason? \}}}\label{checkfreshnessenvelope-now-maxskewms-ok-reason}

Test whether an envelope's signed \texttt{ts} is within the freshness
window (\texttt{MAX\_PUBLISH\_SKEW\_MS}, 5 min by default). Live-publish
ingress enforces this in the kernel; call it yourself only if you need
the check on a path the kernel doesn't gate.

\subsubsection{Envelope constants}\label{envelope-constants}

\begin{itemize}
\tightlist
\item
  \texttt{ENVELOPE\_DOMAIN} ---
  \texttt{\textquotesingle{}axona:pubsub-envelope:v2\textquotesingle{}}
  (signature domain tag).
\item
  \texttt{MAX\_PUBLISH\_SKEW\_MS} --- \texttt{300000} (the freshness
  window).
\end{itemize}

\begin{center}\rule{0.5\linewidth}{0.5pt}\end{center}

\subsection{8. Errors}\label{errors}

All thrown errors inherit from \texttt{AxonaError}:

\begin{Shaded}
\begin{Highlighting}[]
\KeywordTok{class}\NormalTok{ AxonaError }\KeywordTok{extends} \BuiltInTok{Error}\NormalTok{ \{}
  \DataTypeTok{code}\OperatorTok{:}\NormalTok{     string}\OperatorTok{;}   \CommentTok{// stable UPPER\_SNAKE identifier — switch on this, not message}
\NormalTok{  cause}\OperatorTok{?:}   \BuiltInTok{Error}\OperatorTok{;}
\NormalTok{  context}\OperatorTok{?:}\NormalTok{ object}\OperatorTok{;}   \CommentTok{// machine{-}readable details}
\NormalTok{\}}
\end{Highlighting}
\end{Shaded}

Subclasses and their codes:

{\def\LTcaptype{none} % do not increment counter
\begin{longtable}[]{@{}
  >{\raggedright\arraybackslash}p{(\linewidth - 2\tabcolsep) * \real{0.5000}}
  >{\raggedright\arraybackslash}p{(\linewidth - 2\tabcolsep) * \real{0.5000}}@{}}
\toprule\noalign{}
\begin{minipage}[b]{\linewidth}\raggedright
Class
\end{minipage} & \begin{minipage}[b]{\linewidth}\raggedright
Codes
\end{minipage} \\
\midrule\noalign{}
\endhead
\bottomrule\noalign{}
\endlastfoot
\texttt{IdentityError} & \texttt{IDENTITY\_KEYGEN\_FAILED},
\texttt{IDENTITY\_LOAD\_FAILED}, \texttt{IDENTITY\_INVALID\_FORMAT} \\
\texttt{TransportError} & \texttt{TRANSPORT\_NOT\_STARTED},
\texttt{TRANSPORT\_PEER\_UNREACHABLE}, \texttt{TRANSPORT\_TIMEOUT},
\texttt{TRANSPORT\_CHANNEL\_CLOSED},
\texttt{TRANSPORT\_HELLO\_FAILED} \\
\texttt{PublishError} & \texttt{PUBLISH\_INVALID\_TOPIC},
\texttt{PUBLISH\_SIGN\_FAILED}, \texttt{PUBLISH\_NO\_PUBLISH\_IDENTITY},
\texttt{TOPIC\_REGION\_REQUIRED}, \texttt{WRITE\_POLICY\_VIOLATION},
\texttt{PUBLISH\_REPLICATION\_FAILED},
\texttt{PUBLISH\_PAYLOAD\_TOO\_LARGE},
\texttt{PUBLISH\_INVALID\_MESSAGE} \\
\texttt{SubscribeError} & \texttt{SUBSCRIBE\_INVALID\_TOPIC},
\texttt{SUBSCRIBE\_ATTACH\_FAILED},
\texttt{SUBSCRIBE\_HANDLER\_MISSING} \\
\texttt{KillError} & \texttt{KILL\_INVALID\_TOPIC},
\texttt{KILL\_INVALID\_MSGID}, \texttt{KILL\_SIGN\_FAILED} \\
\texttt{UnpubError} & \texttt{UNPUB\_INVALID\_TOPIC},
\texttt{UNPUB\_PUBLIC\_TOPIC}, \texttt{UNPUB\_SIGN\_FAILED} \\
\texttt{TouchError} & \texttt{TOUCH\_INVALID\_TOPIC},
\texttt{TOUCH\_INVALID\_MSGID}, \texttt{TOUCH\_SIGN\_FAILED} \\
\texttt{PullError} & \texttt{PULL\_INVALID\_MSGID},
\texttt{PULL\_AXONS\_UNREACHABLE} \\
\texttt{MetricsError} & \texttt{METRICS\_AXONS\_UNREACHABLE} \\
\texttt{UpgradeRequiredError} & \texttt{UPGRADE\_REQUIRED} (bridge close
code 4426) \\
\end{longtable}
}

The full taxonomy is in \texttt{ErrorCodes}:

\begin{Shaded}
\begin{Highlighting}[]
\ImportTok{import}\NormalTok{ \{ ErrorCodes \} }\ImportTok{from} \StringTok{\textquotesingle{}@axona/protocol\textquotesingle{}}\OperatorTok{;}
\BuiltInTok{console}\OperatorTok{.}\FunctionTok{log}\NormalTok{(ErrorCodes}\OperatorTok{.}\AttributeTok{WRITE\_POLICY\_VIOLATION}\NormalTok{)}\OperatorTok{;}  \CommentTok{// \textquotesingle{}WRITE\_POLICY\_VIOLATION\textquotesingle{}}
\end{Highlighting}
\end{Shaded}

Switch on \texttt{.code}, not \texttt{.message}:

\begin{Shaded}
\begin{Highlighting}[]
\ControlFlowTok{try}\NormalTok{ \{ }\ControlFlowTok{await}\NormalTok{ peer}\OperatorTok{.}\FunctionTok{pub}\NormalTok{(topic}\OperatorTok{,}\NormalTok{ msg}\OperatorTok{,}\NormalTok{ \{ }\DataTypeTok{signWith}\OperatorTok{:}\NormalTok{ me \})}\OperatorTok{;}\NormalTok{ \}}
\ControlFlowTok{catch}\NormalTok{ (err) \{}
  \ControlFlowTok{if}\NormalTok{ (err}\OperatorTok{.}\AttributeTok{code} \OperatorTok{===}\NormalTok{ ErrorCodes}\OperatorTok{.}\AttributeTok{WRITE\_POLICY\_VIOLATION}\NormalTok{) \{ }\CommentTok{/* not the owner */}\NormalTok{ \}}
  \ControlFlowTok{else} \ControlFlowTok{if}\NormalTok{ (err}\OperatorTok{.}\AttributeTok{code} \OperatorTok{===}\NormalTok{ ErrorCodes}\OperatorTok{.}\AttributeTok{PUBLISH\_NO\_PUBLISH\_IDENTITY}\NormalTok{) \{ }\CommentTok{/* name a signer */}\NormalTok{ \}}
  \ControlFlowTok{else} \ControlFlowTok{throw}\NormalTok{ err}\OperatorTok{;}
\NormalTok{\}}
\end{Highlighting}
\end{Shaded}

\subsubsection{Wire-format helpers}\label{wire-format-helpers}

{\def\LTcaptype{none} % do not increment counter
\begin{longtable}[]{@{}
  >{\raggedright\arraybackslash}p{(\linewidth - 2\tabcolsep) * \real{0.5000}}
  >{\raggedright\arraybackslash}p{(\linewidth - 2\tabcolsep) * \real{0.5000}}@{}}
\toprule\noalign{}
\begin{minipage}[b]{\linewidth}\raggedright
Symbol
\end{minipage} & \begin{minipage}[b]{\linewidth}\raggedright
Notes
\end{minipage} \\
\midrule\noalign{}
\endhead
\bottomrule\noalign{}
\endlastfoot
\texttt{isWireError(obj)} → \texttt{boolean} & does this look like an
over-the-wire \texttt{AxonaError}? \\
\texttt{fromWire(obj)} → \texttt{AxonaError} & reconstruct the typed
error (unknown classes fall back to \texttt{AxonaError}). \\
\texttt{err.toWire()} → \texttt{object} &
\texttt{\{\ \_\_axonaError:\ true,\ class,\ code,\ message,\ context\ \}}. \\
\end{longtable}
}

Errors thrown inside a remote \texttt{onMessage} (reached via
\texttt{peer.send}) survive serialization with their class and code
intact.

\begin{center}\rule{0.5\linewidth}{0.5pt}\end{center}

\subsection{9. Persistence}\label{persistence}

\subsubsection{\texorpdfstring{\texttt{PersistenceAdapter} (abstract
class)}{PersistenceAdapter (abstract class)}}\label{persistenceadapter-abstract-class}

The storage contract. Implement it to plug in your own backend.

\begin{Shaded}
\begin{Highlighting}[]
\KeywordTok{class}\NormalTok{ PersistenceAdapter \{}
  \KeywordTok{async} \FunctionTok{load}\NormalTok{(key)        \{ }\CommentTok{/* → previously{-}saved value or null */}\NormalTok{ \}}
  \KeywordTok{async} \FunctionTok{save}\NormalTok{(key}\OperatorTok{,}\NormalTok{ value) \{ }\CommentTok{/* persist value */}\NormalTok{ \}}
  \KeywordTok{async} \KeywordTok{delete}\NormalTok{(key)      \{ }\CommentTok{/* remove key */}\NormalTok{ \}}
\NormalTok{\}}
\end{Highlighting}
\end{Shaded}

Pass an instance as the \texttt{persist} option of the
\texttt{AxonaPeer} constructor; the peer auto-checkpoints identity,
subscriptions, synaptome, and hosting state.

\subsubsection{\texorpdfstring{\texttt{InMemoryPersistence}}{InMemoryPersistence}}\label{inmemorypersistence}

Reference implementation (environment-neutral) used by tests and by the
default \texttt{persist:\ false} path.

\begin{Shaded}
\begin{Highlighting}[]
\ImportTok{import}\NormalTok{ \{ AxonaPeer}\OperatorTok{,}\NormalTok{ InMemoryPersistence \} }\ImportTok{from} \StringTok{\textquotesingle{}@axona/protocol\textquotesingle{}}\OperatorTok{;}
\KeywordTok{const}\NormalTok{ peer }\OperatorTok{=} \KeywordTok{new} \FunctionTok{AxonaPeer}\NormalTok{(\{ nodeIdentity}\OperatorTok{,}\NormalTok{ domain}\OperatorTok{,}\NormalTok{ node}\OperatorTok{,}\NormalTok{ transport}\OperatorTok{,}
  \DataTypeTok{persist}\OperatorTok{:} \KeywordTok{new} \FunctionTok{InMemoryPersistence}\NormalTok{() \})}\OperatorTok{;}
\end{Highlighting}
\end{Shaded}

\subsubsection{\texorpdfstring{\texttt{FilePersistence} /
\texttt{IndexedDBPersistence} (sub-path
imports)}{FilePersistence / IndexedDBPersistence (sub-path imports)}}\label{filepersistence-indexeddbpersistence-sub-path-imports}

Platform-specific implementations --- sub-path imports only (so neither
pulls \texttt{node:fs} into a browser bundle nor
\texttt{globalThis.indexedDB} into Node):

\begin{Shaded}
\begin{Highlighting}[]
\ImportTok{import}\NormalTok{ \{ FilePersistence \}      }\ImportTok{from} \StringTok{\textquotesingle{}@axona/protocol/persistence/file.js\textquotesingle{}}\OperatorTok{;}      \CommentTok{// Node}
\ImportTok{import}\NormalTok{ \{ IndexedDBPersistence \} }\ImportTok{from} \StringTok{\textquotesingle{}@axona/protocol/persistence/indexeddb.js\textquotesingle{}}\OperatorTok{;} \CommentTok{// browser}

\KeywordTok{const}\NormalTok{ peer }\OperatorTok{=} \KeywordTok{new} \FunctionTok{AxonaPeer}\NormalTok{(\{ nodeIdentity}\OperatorTok{,}\NormalTok{ domain}\OperatorTok{,}\NormalTok{ node}\OperatorTok{,}\NormalTok{ transport}\OperatorTok{,}
  \DataTypeTok{persist}\OperatorTok{:} \KeywordTok{new} \FunctionTok{FilePersistence}\NormalTok{(\{ }\DataTypeTok{dir}\OperatorTok{:} \StringTok{\textquotesingle{}./state\textquotesingle{}}\NormalTok{ \}) \})}\OperatorTok{;}
\end{Highlighting}
\end{Shaded}

\begin{center}\rule{0.5\linewidth}{0.5pt}\end{center}

\subsection{Transport + protocol
surface}\label{transport-protocol-surface-1}

\begin{quote}
\textbf{You probably don't need this part.} Everything an application
calls is in the Application surface (§§1--9). These sections are for
people implementing transports, embedding the kernel, or building
tooling.
\end{quote}

\subsection{10. Contracts}\label{contracts}

Abstract base classes that transports / DHT implementations extend.
Application code rarely uses these directly.

\subsubsection{\texorpdfstring{\texttt{Transport} (abstract
class)}{Transport (abstract class)}}\label{transport-abstract-class}

\begin{Shaded}
\begin{Highlighting}[]
\KeywordTok{class}\NormalTok{ MyTransport }\KeywordTok{extends}\NormalTok{ Transport \{}
  \KeywordTok{async} \FunctionTok{start}\NormalTok{(localNodeId)         \{\}}
  \KeywordTok{async} \FunctionTok{stop}\NormalTok{()                     \{\}}
  \KeywordTok{async} \FunctionTok{send}\NormalTok{(peerId}\OperatorTok{,}\NormalTok{ type}\OperatorTok{,}\NormalTok{ body)   \{\}   }\CommentTok{// request/response}
  \KeywordTok{async} \FunctionTok{notify}\NormalTok{(peerId}\OperatorTok{,}\NormalTok{ type}\OperatorTok{,}\NormalTok{ body) \{\}   }\CommentTok{// fire{-}and{-}forget}
  \KeywordTok{async} \FunctionTok{openConnection}\NormalTok{(peerId)     \{\}}
  \KeywordTok{async} \FunctionTok{closeConnection}\NormalTok{(peerId)    \{\}}
  \FunctionTok{isConnected}\NormalTok{(peerId)              \{\}}
  \FunctionTok{getLatency}\NormalTok{(peerId)               \{\}   }\CommentTok{// RTT ms or {-}1}
  \FunctionTok{onRequest}\NormalTok{(type}\OperatorTok{,}\NormalTok{ handler)         \{\}}
  \FunctionTok{onNotification}\NormalTok{(type}\OperatorTok{,}\NormalTok{ handler)    \{\}}
  \FunctionTok{onPeerDied}\NormalTok{(handler)              \{\}}
\NormalTok{\}}
\end{Highlighting}
\end{Shaded}

Full surface in \texttt{@axona/protocol/contracts/Transport.js}.

\subsubsection{\texorpdfstring{\texttt{DHT} (abstract
class)}{DHT (abstract class)}}\label{dht-abstract-class}

The per-node routing/pub-sub contract \texttt{AxonaPeer} implements
(\texttt{getSelfId}, \texttt{findKClosest}, \texttt{routeMessage},
\texttt{sendDirect}, \texttt{onRoutedMessage}, \texttt{lookup}, \ldots).

\subsubsection{\texorpdfstring{\texttt{BootstrapService} (abstract
class)}{BootstrapService (abstract class)}}\label{bootstrapservice-abstract-class}

For peer discovery (signaling brokers, seed lists, DNS-SD). Rarely
needed by app code.

\begin{center}\rule{0.5\linewidth}{0.5pt}\end{center}

\subsection{11. Sim transport}\label{sim-transport}

In-process router for tests + simulators. Environment-neutral, ships in
the main barrel.

\subsubsection{\texorpdfstring{\texttt{new\ SimNetwork(\{\ latencyFn?\ \})}}{new SimNetwork(\{ latencyFn? \})}}\label{new-simnetwork-latencyfn}

\begin{Shaded}
\begin{Highlighting}[]
\ImportTok{import}\NormalTok{ \{ SimNetwork \} }\ImportTok{from} \StringTok{\textquotesingle{}@axona/protocol\textquotesingle{}}\OperatorTok{;}
\KeywordTok{const}\NormalTok{ net }\OperatorTok{=} \KeywordTok{new} \FunctionTok{SimNetwork}\NormalTok{()}\OperatorTok{;}                          \CommentTok{// 0 ms edges}
\KeywordTok{const}\NormalTok{ net2 }\OperatorTok{=} \KeywordTok{new} \FunctionTok{SimNetwork}\NormalTok{(\{ }\DataTypeTok{latencyFn}\OperatorTok{:}\NormalTok{ () }\KeywordTok{=\textgreater{}} \DecValTok{50}\NormalTok{ \})}\OperatorTok{;}  \CommentTok{// 50 ms each way}
\end{Highlighting}
\end{Shaded}

\subsubsection{\texorpdfstring{\texttt{simTransport(\{\ network,\ identity,\ heartbeatMs?,\ ...\ \})}
→
\texttt{SimTransport}}{simTransport(\{ network, identity, heartbeatMs?, ... \}) → SimTransport}}\label{simtransport-network-identity-heartbeatms-...-simtransport}

Factory that builds + wires a \texttt{SimTransport} onto a
\texttt{SimNetwork}.

\begin{Shaded}
\begin{Highlighting}[]
\ImportTok{import}\NormalTok{ \{ SimNetwork}\OperatorTok{,}\NormalTok{ simTransport \} }\ImportTok{from} \StringTok{\textquotesingle{}@axona/protocol\textquotesingle{}}\OperatorTok{;}
\KeywordTok{const}\NormalTok{ network }\OperatorTok{=} \KeywordTok{new} \FunctionTok{SimNetwork}\NormalTok{()}\OperatorTok{;}
\KeywordTok{const}\NormalTok{ t }\OperatorTok{=} \FunctionTok{simTransport}\NormalTok{(\{ network}\OperatorTok{,} \DataTypeTok{identity}\OperatorTok{:}\NormalTok{ node}\OperatorTok{,} \DataTypeTok{heartbeatMs}\OperatorTok{:} \DecValTok{0}\NormalTok{ \})}\OperatorTok{;}
\ControlFlowTok{await}\NormalTok{ t}\OperatorTok{.}\FunctionTok{start}\NormalTok{(node}\OperatorTok{.}\AttributeTok{id}\NormalTok{)}\OperatorTok{;}
\ControlFlowTok{await}\NormalTok{ t}\OperatorTok{.}\FunctionTok{openConnection}\NormalTok{(other}\OperatorTok{.}\AttributeTok{id}\NormalTok{)}\OperatorTok{;}
\end{Highlighting}
\end{Shaded}

\texttt{heartbeatMs:\ 0} disables heartbeats (recommended for tests).
\texttt{SimTransport} is also exported directly for
\texttt{new\ SimTransport(\{...\})}.

\begin{center}\rule{0.5\linewidth}{0.5pt}\end{center}

\subsection{12. Web transport}\label{web-transport}

The production browser transport (a WebSocket to the bridge for
bootstrap + signaling, plus a WebRTC mesh). Sub-path import:

\begin{Shaded}
\begin{Highlighting}[]
\ImportTok{import}\NormalTok{ \{ webTransport \} }\ImportTok{from} \StringTok{\textquotesingle{}@axona/protocol/transport/web/index.js\textquotesingle{}}\OperatorTok{;}

\KeywordTok{const}\NormalTok{ transport }\OperatorTok{=} \FunctionTok{webTransport}\NormalTok{(\{}
  \DataTypeTok{bridgeUrl}\OperatorTok{:}   \StringTok{\textquotesingle{}wss://testnet.axona.net\textquotesingle{}}\OperatorTok{,}  \CommentTok{// or wss://bridge.axona.net (production)}
  \DataTypeTok{identity}\OperatorTok{:}\NormalTok{    node}\OperatorTok{,}                       \CommentTok{// from createNodeIdentity — signs the handshake}
  \DataTypeTok{peerVersion}\OperatorTok{:} \StringTok{\textquotesingle{}4.22.0\textquotesingle{}}\OperatorTok{,}                   \CommentTok{// your app version (gated by the bridge)}
  \DataTypeTok{reconnect}\OperatorTok{:}   \KeywordTok{true}\OperatorTok{,}
\NormalTok{\})}\OperatorTok{;}
\ControlFlowTok{await}\NormalTok{ transport}\OperatorTok{.}\FunctionTok{start}\NormalTok{()}\OperatorTok{;}                  \CommentTok{// resolves after the bridge handshake}
\end{Highlighting}
\end{Shaded}

Beyond the \texttt{Transport} contract it exposes an observability
surface: \texttt{transport.bridgeState}, \texttt{bridgeInfo},
\texttt{bridgeRtt}, \texttt{bridgeNodeId}, \texttt{onBridgeState(cb)},
\texttt{onWelcome(cb)}, \texttt{boundPeers()}, \texttt{reconnectNow()}.

Each mesh link's \texttt{axona/5} proof is bound to its DTLS certificate
fingerprint, so a fingerprint-rewriting bridge cannot transparently MITM
``direct'' peer traffic --- the bridge is trusted for signaling only.

\begin{center}\rule{0.5\linewidth}{0.5pt}\end{center}

\subsection{13. Handshake}\label{handshake}

Wire-version handshake helpers used by transports on a fresh signaling
channel. Most app code never touches these --- they're plumbed into the
transport factories.

\begin{Shaded}
\begin{Highlighting}[]
\FunctionTok{buildClientHello}\NormalTok{(\{ version}\OperatorTok{,}\NormalTok{ wireVersion}\OperatorTok{?,}\NormalTok{ capabilities}\OperatorTok{?}\NormalTok{ \})          }\CommentTok{// → client → server frame}
\FunctionTok{buildServerHello}\NormalTok{(\{ version}\OperatorTok{,}\NormalTok{ wireVersion}\OperatorTok{?,}\NormalTok{ minPeerVersion}\OperatorTok{,}\NormalTok{ downloadUrl \}) }\CommentTok{// → server → client frame}
\FunctionTok{parseHello}\NormalTok{(frame)                                                   }\CommentTok{// → parsed hello}
\FunctionTok{parseVersion}\NormalTok{(str)                                                   }\CommentTok{// → \{ major, minor, patch \}}
\FunctionTok{compareVersions}\NormalTok{(a}\OperatorTok{,}\NormalTok{ b)                                               }\CommentTok{// → {-}1 | 0 | 1}
\FunctionTok{wireCompatible}\NormalTok{(a}\OperatorTok{,}\NormalTok{ b)                                                }\CommentTok{// → boolean}
\FunctionTok{performClientHandshake}\NormalTok{(channel}\OperatorTok{,}\NormalTok{ opts)                               }\CommentTok{// → handshake result}
\FunctionTok{performServerHandshake}\NormalTok{(channel}\OperatorTok{,}\NormalTok{ opts)}
\end{Highlighting}
\end{Shaded}

The constants \texttt{WIRE\_VERSION}, \texttt{KERNEL\_VERSION},
\texttt{UPGRADE\_CLOSE\_CODE} come from the same module --- see
\hyperref[23-constants]{§23}.

\begin{center}\rule{0.5\linewidth}{0.5pt}\end{center}

\subsection{14. Authenticated-identity handshake
(axona/5)}\label{authenticated-identity-handshake-axona5}

Re-exported so consumers driving their own channel lifecycle (the
bridge's embedded peer, custom transports) can build/verify
authenticated hellos with the same primitive the web transport uses.

\begin{Shaded}
\begin{Highlighting}[]
\FunctionTok{buildAuthHello}\NormalTok{(opts)                  }\CommentTok{// → signed hello proving nodeId ownership}
\FunctionTok{verifyAuthHello}\NormalTok{(hello}\OperatorTok{,}\NormalTok{ opts)          }\CommentTok{// → verification result}
\FunctionTok{pubkeyMatchesNodeId}\NormalTok{(pubkey}\OperatorTok{,}\NormalTok{ nodeId)   }\CommentTok{// → boolean (BIND check)}
\FunctionTok{makeNonce}\NormalTok{()                           }\CommentTok{// → fresh challenge nonce}
\FunctionTok{cbvFromNonces}\NormalTok{(}\OperatorTok{...}\NormalTok{)                    }\CommentTok{// → channel{-}binding value from nonces}
\FunctionTok{cbvFromFingerprints}\NormalTok{(}\OperatorTok{...}\NormalTok{)              }\CommentTok{// → channel{-}binding value from DTLS fingerprints}
\NormalTok{AUTH\_PROTO                            }\CommentTok{// \textquotesingle{}axona/5\textquotesingle{}}
\end{Highlighting}
\end{Shaded}

The proof binds (BIND: pubkey hashes to the id), (POSSESS: Ed25519
signature), and (CHANNEL: bound to this live connection) so a captured
proof can't be replayed onto another link.

\begin{center}\rule{0.5\linewidth}{0.5pt}\end{center}

\subsection{15. Bridge directory}\label{bridge-directory}

Bridges advertise on a public topic; clients collect them at launch and
fail over.

\begin{Shaded}
\begin{Highlighting}[]
\NormalTok{BRIDGE\_DIRECTORY\_TOPIC               }\CommentTok{// \textquotesingle{}axona:bridge{-}directory\textquotesingle{}}
\NormalTok{BRIDGE\_ENTRY\_MAX\_AGE\_MS              }\CommentTok{// 48 h}
\FunctionTok{buildBridgeEntry}\NormalTok{(\{ url}\OperatorTok{,}\NormalTok{ lat}\OperatorTok{,}\NormalTok{ lng}\OperatorTok{,}\NormalTok{ label}\OperatorTok{?,}\NormalTok{ ver}\OperatorTok{?,}\NormalTok{ turn}\OperatorTok{?,}\NormalTok{ ts}\OperatorTok{?}\NormalTok{ \})  }\CommentTok{// → directory entry}
\FunctionTok{validateBridgeEntry}\NormalTok{(msg)             }\CommentTok{// → boolean (well{-}formed + fresh)}
\FunctionTok{rankBridges}\NormalTok{(\{ }\OperatorTok{...}\NormalTok{ \})                 }\CommentTok{// → ordered candidate list (reputation + proximity)}
\FunctionTok{haversineKm}\NormalTok{(a}\OperatorTok{,}\NormalTok{ b)                    }\CommentTok{// → great{-}circle km between two \{ lat, lng \}}
\end{Highlighting}
\end{Shaded}

\begin{Shaded}
\begin{Highlighting}[]
\ImportTok{import}\NormalTok{ \{ BRIDGE\_DIRECTORY\_TOPIC}\OperatorTok{,}\NormalTok{ buildBridgeEntry \} }\ImportTok{from} \StringTok{\textquotesingle{}@axona/protocol\textquotesingle{}}\OperatorTok{;}
\KeywordTok{const}\NormalTok{ entry }\OperatorTok{=} \FunctionTok{buildBridgeEntry}\NormalTok{(\{ }\DataTypeTok{url}\OperatorTok{:} \StringTok{\textquotesingle{}wss://bridge.axona.net\textquotesingle{}}\OperatorTok{,} \DataTypeTok{lat}\OperatorTok{:} \DecValTok{38}\OperatorTok{,} \DataTypeTok{lng}\OperatorTok{:} \OperatorTok{{-}}\DecValTok{77}\NormalTok{ \})}\OperatorTok{;}
\ControlFlowTok{await}\NormalTok{ peer}\OperatorTok{.}\FunctionTok{pub}\NormalTok{(\{ }\DataTypeTok{name}\OperatorTok{:}\NormalTok{ BRIDGE\_DIRECTORY\_TOPIC}\OperatorTok{,} \DataTypeTok{region}\OperatorTok{:} \StringTok{\textquotesingle{}useast\textquotesingle{}}\NormalTok{ \}}\OperatorTok{,}\NormalTok{ entry}\OperatorTok{,}\NormalTok{ \{ }\DataTypeTok{signWith}\OperatorTok{:}\NormalTok{ me \})}\OperatorTok{;}
\end{Highlighting}
\end{Shaded}

\begin{center}\rule{0.5\linewidth}{0.5pt}\end{center}

\subsection{16. Low-level pub/sub
builders}\label{low-level-pubsub-builders}

The signed wire record behind \texttt{peer.kill}. Apps use
\texttt{peer.kill}; this is for protocol implementors and custom roots.

\begin{Shaded}
\begin{Highlighting}[]
\FunctionTok{buildKill}\NormalTok{(\{ topicId}\OperatorTok{,}\NormalTok{ msgId}\OperatorTok{,}\NormalTok{ ts}\OperatorTok{?,}\NormalTok{ seq}\OperatorTok{?,}\NormalTok{ identity \})                       }\CommentTok{// → signed kill}
\FunctionTok{verifyKill}\NormalTok{(kill)                                                          }\CommentTok{// → verification result}
\NormalTok{KILL\_DOMAIN     }\CommentTok{// \textquotesingle{}axona:pubsub{-}kill:v1\textquotesingle{}}
\end{Highlighting}
\end{Shaded}

\begin{quote}
The \texttt{unpub} wire record was \textbf{removed in v4.3.0} (kill is
the single retraction primitive). \texttt{touch} is \textbf{deprecated}
--- its handler is a no-op; the \texttt{buildTouch}/\texttt{verifyTouch}
helpers remain only for wire back-compat.
\end{quote}

\texttt{AxonaManager} (the pub/sub state machine), \texttt{AxonaDomain},
\texttt{NeuronNode}, \texttt{Synapse}, \texttt{Subscription},
\texttt{DHTNode}, and \texttt{GEO\_CELL\_BITS} (= 8) are also exported
for implementors building custom engines.

\begin{center}\rule{0.5\linewidth}{0.5pt}\end{center}

\subsection{Low-level utilities}\label{low-level-utilities-1}

\begin{quote}
\textbf{Internals.} Pure helpers the kernel itself is built from. Apps
import these only for unusual jobs (custom signing, ID math, region
tables).
\end{quote}

\subsection{17. Ed25519 helpers}\label{ed25519-helpers}

Web Crypto Ed25519 wrappers --- useful for signing/verifying outside
\texttt{buildEnvelope}.

\begin{Shaded}
\begin{Highlighting}[]
\KeywordTok{const}\NormalTok{ \{ publicKey}\OperatorTok{,}\NormalTok{ privateKey \} }\OperatorTok{=} \ControlFlowTok{await} \FunctionTok{generateKeyPair}\NormalTok{(\{ }\DataTypeTok{extractable}\OperatorTok{:} \KeywordTok{true}\NormalTok{ \})}\OperatorTok{;}
\KeywordTok{const}\NormalTok{ pubHex  }\OperatorTok{=} \ControlFlowTok{await} \FunctionTok{exportPublicKey}\NormalTok{(publicKey)}\OperatorTok{;}   \CommentTok{// 64{-}char hex}
\KeywordTok{const}\NormalTok{ pubKey  }\OperatorTok{=} \ControlFlowTok{await} \FunctionTok{importPublicKey}\NormalTok{(pubBytesOrHex)}\OperatorTok{;}
\KeywordTok{const}\NormalTok{ sig     }\OperatorTok{=} \ControlFlowTok{await} \FunctionTok{sign}\NormalTok{(privateKey}\OperatorTok{,}\NormalTok{ dataBytes)}\OperatorTok{;}  \CommentTok{// Uint8Array(64)}
\KeywordTok{const}\NormalTok{ ok      }\OperatorTok{=} \ControlFlowTok{await} \FunctionTok{verify}\NormalTok{(pubKeyOrBytes}\OperatorTok{,}\NormalTok{ dataBytes}\OperatorTok{,}\NormalTok{ sig)}\OperatorTok{;}

\KeywordTok{const}\NormalTok{ signer   }\OperatorTok{=} \FunctionTok{makeSigner}\NormalTok{(privateKey)}\OperatorTok{;}   \CommentTok{// (bytes) =\textgreater{} Promise\textless{}sig\textgreater{}}
\KeywordTok{const}\NormalTok{ verifier }\OperatorTok{=} \FunctionTok{makeVerifier}\NormalTok{(pubKey)}\OperatorTok{;}      \CommentTok{// (bytes, sig) =\textgreater{} Promise\textless{}boolean\textgreater{}}
\end{Highlighting}
\end{Shaded}

Implementation-agnostic: substitute \texttt{@noble/ed25519} on runtimes
without Web Crypto Ed25519 (Chrome \textless110, Safari \textless17,
Firefox \textless130, Node \textless20).

\begin{center}\rule{0.5\linewidth}{0.5pt}\end{center}

\subsection{18. Hex / ID math}\label{hex-id-math}

264-bit identifier math --- nodeId and topicId share the same keyspace
(\texttt{{[}8-bit\ S2\ prefix{]}\ \textbar{}\textbar{}\ {[}256-bit\ hash{]}}).

\begin{Shaded}
\begin{Highlighting}[]
\NormalTok{ID\_BITS                     }\CommentTok{// 264}
\NormalTok{HASH\_BITS                   }\CommentTok{// 256}
\NormalTok{S2\_BITS                     }\CommentTok{// 8}
\NormalTok{HEX\_CHARS                   }\CommentTok{// 66}
\NormalTok{MAX\_ID                      }\CommentTok{// BigInt (2**264 {-} 1)}
\NormalTok{MAX\_HASH                    }\CommentTok{// BigInt (2**256 {-} 1)}
\NormalTok{MAX\_S2                      }\CommentTok{// 255}

\FunctionTok{toHex}\NormalTok{(bigint)               }\CommentTok{// → 66{-}char lowercase hex}
\FunctionTok{fromHex}\NormalTok{(hex)                }\CommentTok{// → BigInt}
\FunctionTok{isHexId}\NormalTok{(s)                  }\CommentTok{// → boolean (valid 66{-}char hex id)}
\FunctionTok{assembleId}\NormalTok{(s2Prefix}\OperatorTok{,}\NormalTok{ hash)  }\CommentTok{// → BigInt — 8{-}bit prefix || 256{-}bit hash}
\FunctionTok{extractS2Prefix}\NormalTok{(idBig)      }\CommentTok{// → number 0..255 (top byte)}
\FunctionTok{extractHash}\NormalTok{(idBig)          }\CommentTok{// → BigInt (bottom 256 bits)}
\FunctionTok{s2PrefixOfHex}\NormalTok{(hex)          }\CommentTok{// → number 0..255}
\FunctionTok{xorDistance}\NormalTok{(a}\OperatorTok{,}\NormalTok{ b)           }\CommentTok{// → BigInt}
\FunctionTok{stratumOf}\NormalTok{(a}\OperatorTok{,}\NormalTok{ b)             }\CommentTok{// → 0..264 — leading{-}zero count of (a \^{} b)}
\FunctionTok{clz264}\NormalTok{(bigint)              }\CommentTok{// → 0..264 — leading{-}zero count}
\FunctionTok{randomU256}\NormalTok{()                }\CommentTok{// → BigInt — cryptographically random}
\end{Highlighting}
\end{Shaded}

\begin{center}\rule{0.5\linewidth}{0.5pt}\end{center}

\subsection{19. S2 geographic cells}\label{s2-geographic-cells}

The 8-bit S2 cell that occupies the top byte of every id. The cellId
matches the top 8 bits of Google S2's level-3 cell ID (full interop).

\begin{Shaded}
\begin{Highlighting}[]
\FunctionTok{geoCellId}\NormalTok{(lat}\OperatorTok{,}\NormalTok{ lng}\OperatorTok{,}\NormalTok{ bits)   }\CommentTok{// → integer — S2 cell ID at the given bit depth}
\FunctionTok{geoCellCenter}\NormalTok{(cellId)       }\CommentTok{// → \{ lat, lng \} — cell center}
\FunctionTok{geoCellCorners}\NormalTok{(cellId)      }\CommentTok{// → corner coordinates}
\FunctionTok{geoCellFace}\NormalTok{(cellId)         }\CommentTok{// → 0..5 — which cube face}
\FunctionTok{geoCellSubCenters}\NormalTok{(cellId)   }\CommentTok{// → [center0, center1] — the two level{-}3 sub{-}cell centers}
\FunctionTok{geoCellHalf}\NormalTok{(lat}\OperatorTok{,}\NormalTok{ lng)       }\CommentTok{// → 0 | 1 — which sub{-}cell a point falls in}
\FunctionTok{isValidCellId}\NormalTok{(cellId)       }\CommentTok{// → boolean (in [0, 192))}
\NormalTok{S2\_FACES                    }\CommentTok{// 6}
\NormalTok{S2\_CELL\_COUNT               }\CommentTok{// 192}
\NormalTok{S2\_RESERVED\_FROM            }\CommentTok{// 192}
\end{Highlighting}
\end{Shaded}

\begin{center}\rule{0.5\linewidth}{0.5pt}\end{center}

\subsection{20. Region names}\label{region-names}

Each of the 192 region codes carries exactly one human-readable name, so
a region always presents the same label. Where a cell straddles ocean
and land the land name wins; a multi-country cell takes its dominant
city.

\begin{Shaded}
\begin{Highlighting}[]
\NormalTok{REGION\_NAMES                  }\CommentTok{// frozen string[] indexed by code [0,192)}
\FunctionTok{regionName}\NormalTok{(code)              }\CommentTok{// → string | null     e.g. regionName(0x89) → \textquotesingle{}useast\textquotesingle{}}
\FunctionTok{regionNames}\NormalTok{(code)             }\CommentTok{// → [name] | null      deprecated one{-}element shim}
\FunctionTok{regionCode}\NormalTok{(name)              }\CommentTok{// → number | null      inverse (canonical code for a multi{-}cell name)}
\FunctionTok{resolveRegion}\NormalTok{(nameOrCode)     }\CommentTok{// → number | null      accepts \textquotesingle{}useast\textquotesingle{} | \textquotesingle{}0x89\textquotesingle{} | \textquotesingle{}137\textquotesingle{} | 137}
\FunctionTok{regionNameForLatLng}\NormalTok{(lat}\OperatorTok{,}\NormalTok{ lng) }\CommentTok{// → string             region name for a coordinate}
\FunctionTok{regionCenter}\NormalTok{(nameOrCode)      }\CommentTok{// → \{ lat, lng \} | null cell center (default placement for topics)}
\NormalTok{POPULATED\_REGIONS             }\CommentTok{// frozen [\{ code, name \}] for every non{-}open{-}ocean cell}
\end{Highlighting}
\end{Shaded}

\begin{Shaded}
\begin{Highlighting}[]
\ImportTok{import}\NormalTok{ \{ regionCode}\OperatorTok{,}\NormalTok{ regionCenter}\OperatorTok{,}\NormalTok{ POPULATED\_REGIONS \} }\ImportTok{from} \StringTok{\textquotesingle{}@axona/protocol\textquotesingle{}}\OperatorTok{;}
\FunctionTok{regionCode}\NormalTok{(}\StringTok{\textquotesingle{}useast\textquotesingle{}}\NormalTok{)}\OperatorTok{;}         \CommentTok{// 137}
\FunctionTok{regionCenter}\NormalTok{(}\StringTok{\textquotesingle{}useast\textquotesingle{}}\NormalTok{)}\OperatorTok{;}       \CommentTok{// \{ lat, lng \} — mint a node identity in a region}
\NormalTok{POPULATED\_REGIONS}\OperatorTok{.}\AttributeTok{length}\OperatorTok{;}     \CommentTok{// count of real, inhabited cells (for a region picker)}
\end{Highlighting}
\end{Shaded}

Names match \texttt{/\^{}{[}a-z0-9\_{]}\{1,8\}\$/}; open-ocean cells are
\texttt{\textless{}ocean3\textgreater{}\_\textless{}hex\textgreater{}}
(\texttt{pac\_68}, \texttt{atl\_0a}, \ldots) and are excluded from
\texttt{POPULATED\_REGIONS}.

\begin{center}\rule{0.5\linewidth}{0.5pt}\end{center}

\subsection{21. Geo helpers}\label{geo-helpers}

\begin{Shaded}
\begin{Highlighting}[]
\FunctionTok{haversine}\NormalTok{(lat1}\OperatorTok{,}\NormalTok{ lng1}\OperatorTok{,}\NormalTok{ lat2}\OperatorTok{,}\NormalTok{ lng2)        }\CommentTok{// → km between two points}
\FunctionTok{roundTripLatency}\NormalTok{(lat1}\OperatorTok{,}\NormalTok{ lng1}\OperatorTok{,}\NormalTok{ lat2}\OperatorTok{,}\NormalTok{ lng2) }\CommentTok{// → ms — rough fiber estimate}
\FunctionTok{randomU32}\NormalTok{()                               }\CommentTok{// → 0..2**32 {-} 1}
\end{Highlighting}
\end{Shaded}

The remaining \texttt{geo.js} helpers (continent detection, XOR
routing-table builders, etc.) are reachable via
\texttt{@axona/protocol/utils/geo.js}.

\begin{center}\rule{0.5\linewidth}{0.5pt}\end{center}

\subsection{22. Proof-of-work
scaffolding}\label{proof-of-work-scaffolding}

E-1 / Stage 2 transport + publish PoW. \textbf{Inert at difficulty 0}
--- every nonce is \texttt{\textquotesingle{}\textquotesingle{}} and
verification passes trivially until difficulty is raised. Apps don't
normally touch these; identities carry their PoW nonce automatically.

\begin{Shaded}
\begin{Highlighting}[]
\NormalTok{POW\_DOMAIN                                   }\CommentTok{// \textquotesingle{}axona:pow:v1\textquotesingle{}}
\NormalTok{POW\_DIFFICULTY                               }\CommentTok{// \{ transport: 0, publish: 0 \} (frozen defaults)}
\FunctionTok{powDifficulty}\NormalTok{(role)                          }\CommentTok{// → current difficulty for \textquotesingle{}transport\textquotesingle{} | \textquotesingle{}publish\textquotesingle{}}
\FunctionTok{setPowDifficulty}\NormalTok{(role}\OperatorTok{,}\NormalTok{ n)                    }\CommentTok{// set difficulty (testing / calibration)}
\FunctionTok{resetPowDifficulty}\NormalTok{()                         }\CommentTok{// restore defaults}
\FunctionTok{powMint}\NormalTok{(\{ pubkeyHex}\OperatorTok{,}\NormalTok{ role}\OperatorTok{?,}\NormalTok{ difficulty}\OperatorTok{?,}\NormalTok{ maxTries}\OperatorTok{?}\NormalTok{ \})  }\CommentTok{// → Promise\textless{}nonce string\textgreater{}}
\FunctionTok{powVerify}\NormalTok{(\{ pubkeyHex}\OperatorTok{,}\NormalTok{ nonce}\OperatorTok{,}\NormalTok{ role}\OperatorTok{?,}\NormalTok{ difficulty}\OperatorTok{?}\NormalTok{ \})    }\CommentTok{// → Promise\textless{}boolean\textgreater{}}
\FunctionTok{powBits}\NormalTok{(\{ pubkeyHex}\OperatorTok{,}\NormalTok{ nonce}\OperatorTok{,}\NormalTok{ role}\OperatorTok{?}\NormalTok{ \})         }\CommentTok{// → Promise\textless{}number\textgreater{} leading{-}zero bits of the puzzle hash}
\FunctionTok{powCalibrate}\NormalTok{(\{ ms}\OperatorTok{?}\NormalTok{ \})                        }\CommentTok{// → Promise\textless{}calibration\textgreater{} bits/ms estimate}
\end{Highlighting}
\end{Shaded}

The puzzle hashes
\texttt{H(domain\ \textbar{}\textbar{}\ role\ \textbar{}\textbar{}\ pubkey\ \textbar{}\textbar{}\ nonce)}
--- difficulty is on a \textbf{separate puzzle hash, never on the node
address}, so raising it introduces no keyspace skew.

\begin{center}\rule{0.5\linewidth}{0.5pt}\end{center}

\subsection{23. Constants}\label{constants}

\begin{Shaded}
\begin{Highlighting}[]
\NormalTok{WIRE\_VERSION         }\CommentTok{// \textquotesingle{}4.0\textquotesingle{}      — wire format major.minor (bridges gate on this)}
\NormalTok{KERNEL\_VERSION       }\CommentTok{// \textquotesingle{}4.22.0\textquotesingle{}   — kernel semver (npm release tag)}
\NormalTok{AUTH\_PROTO           }\CommentTok{// \textquotesingle{}axona/5\textquotesingle{}  — authenticated{-}identity handshake tag}
\NormalTok{UPGRADE\_CLOSE\_CODE   }\CommentTok{// 4426       — WebSocket close code for a version mismatch}
\NormalTok{ENVELOPE\_DOMAIN      }\CommentTok{// \textquotesingle{}axona:pubsub{-}envelope:v2\textquotesingle{}}
\NormalTok{MAX\_PUBLISH\_SKEW\_MS  }\CommentTok{// 300000     — freshness window (5 min)}
\NormalTok{ID\_BITS              }\CommentTok{// 264}
\NormalTok{HEX\_CHARS            }\CommentTok{// 66}
\end{Highlighting}
\end{Shaded}

\texttt{WIRE\_VERSION} is what bridges enforce gates against;
\texttt{KERNEL\_VERSION} is informational.

\subsubsection{\texorpdfstring{\texttt{ANONYMOUS}
(sentinel)}{ANONYMOUS (sentinel)}}\label{anonymous-sentinel}

The sentinel for an intentionally unsigned publish. Anonymity must be
explicit --- omitting a signer is an error, never silent anonymity.

\begin{Shaded}
\begin{Highlighting}[]
\ImportTok{import}\NormalTok{ \{ ANONYMOUS \} }\ImportTok{from} \StringTok{\textquotesingle{}@axona/protocol\textquotesingle{}}\OperatorTok{;}
\ControlFlowTok{await}\NormalTok{ peer}\OperatorTok{.}\FunctionTok{pub}\NormalTok{(\{ }\DataTypeTok{region}\OperatorTok{:} \StringTok{\textquotesingle{}useast\textquotesingle{}}\OperatorTok{,} \DataTypeTok{name}\OperatorTok{:} \StringTok{\textquotesingle{}lobby\textquotesingle{}}\NormalTok{ \}}\OperatorTok{,}\NormalTok{ msg}\OperatorTok{,}\NormalTok{ \{ }\DataTypeTok{signWith}\OperatorTok{:}\NormalTok{ ANONYMOUS \})}\OperatorTok{;}
\end{Highlighting}
\end{Shaded}

\begin{center}\rule{0.5\linewidth}{0.5pt}\end{center}

\subsection{Appendix: version history for
migrators}\label{appendix-version-history-for-migrators}

\emph{(Only relevant if you have code from an earlier kernel line. New
readers: skip.)}

\subsubsection{What changed in v4.17.0--v4.22.0 (2026-07-14 roll;
v4.21.0 promoted to
production)}\label{what-changed-in-v4.17.0v4.22.0-2026-07-14-roll-v4.21.0-promoted-to-production}

\textbf{No API change at all} --- every release in this span is
behavioral: root lifecycle, departure, and history-convergence work,
validated by an overnight production-shaped soak per release.

\begin{itemize}
\tightlist
\item
  \textbf{Root election + reconciliation (v4.17.0--v4.19.2).} Faster
  promotion when a topic's root churns out; wrong root claims converge
  without flapping (strictly-closer beacon deferral, periodic root
  self-verification, and guards for unmeshed and cross-region edge
  cases). One durable root per topic, kept true under churn.
\item
  \textbf{Departure-side durability (v4.19.4--v4.19.5).}
  \texttt{peer.leave()} now hands off \textbf{every} rooted topic's
  history inside its time bound (parallel heir resolution + an
  iterative-lookup fallback for thin-tabled leavers), and the heir can
  no longer defer its claim back to the node that just left. A burst
  publisher's topics survive its departure intact.
\item
  \textbf{Kernel consolidation (v4.19.6--v4.21.0).} The refactor
  program: a frozen invariants contract (\texttt{INVARIANTS.md}), every
  root transition through one state machine (\texttt{rootClaim.js}), and
  the pub/sub manager split along its seams. Structural only ---
  verified behavior-preserving.
\item
  \textbf{Split-history union (v4.22.0).} After a root transition, the
  old and new holders converge to the \textbf{union} of cache +
  tombstones (the subscribe now advertises its oldest stamp so a root
  pulls history that sits \emph{below} its own high-water; roots
  union-ingest cohort pushes). Closes the last known replay gap: a fresh
  \texttt{since:\textquotesingle{}all\textquotesingle{}} subscriber
  attaching mid-transition received only the post-transition half
  (\textasciitilde1 in 15 root transitions); soak-validated to
  full-timeline recovery.
\end{itemize}

\subsubsection{What changed in v4.12.0--v4.16.1 (2026-07-02 testnet
roll)}\label{what-changed-in-v4.12.0v4.16.1-2026-07-02-testnet-roll}

\begin{itemize}
\tightlist
\item
  \textbf{\texttt{connect()} (v4.16.0)} --- the one-call bootstrap
  above; no other API change.
\item
  \textbf{Immediate metrics answer (v4.16.1)} --- a topic root publishes
  the FIRST metric snapshot the moment a metrics lease arms, so it
  arrives at routing latency (measured \textasciitilde0.3 s on testnet)
  instead of on the next 5 s tick. Cadence (\textasciitilde20 s) and
  lease (\textasciitilde70 s) unchanged.
\item
  \textbf{Demand-driven metrics (v4.12.0)} --- snapshots publish to
  \texttt{metricTopic(T)} only while a subscriber's lease is fresh;
  \texttt{peer.metrics()} unchanged. The relay's old push-metrics loop
  is retired. See the timing contract in §4.3.
\item
  \textbf{Region-occupancy rule (v4.13.0, gated OFF in v4.15.0)} --- the
  kernel can enforce region-homogeneous topic service
  (\texttt{configureRegionLock(\{\ enforce\ \})}); disabled by default
  until regional coverage justifies it.
\item
  \textbf{BigInt id invariant (v4.14.0)} --- internal ids validated at
  one gate (\texttt{asId}); \texttt{NeuronNode}/\texttt{Synapse} now
  accept hex ids directly. \texttt{asId} exported.
\end{itemize}

\subsubsection{What changed in the 4.x line (since
v3.6.0)}\label{what-changed-in-the-4.x-line-since-v3.6.0}

The \textbf{public API surface is unchanged} from v3.6.0 --- same
identity factories, topic descriptors, and
\texttt{pub}/\texttt{sub}/\texttt{pull}/\texttt{kill}/\texttt{host}
signatures. The 4.x work is a routing-and-reliability rewrite
\emph{under} that surface, so existing code keeps working while delivery
gets more robust:

\begin{itemize}
\tightlist
\item
  \textbf{Routing-only axonic-tree pub/sub (v3.14, clean break).} Topics
  route to an emergent root (the live node XOR-closest to the topic id)
  with no separate overlay to maintain. Behavioral only --- no call
  changed.
\item
  \textbf{K-closest cohort distribution (v4.10.0).} A topic's
  authoritative state (messages \emph{and} retractions) is replicated
  across the closest-K nodes, not one root. A \texttt{kill} reaches the
  whole cohort and internal transfers carry their tombstones, so a
  killed message can't resurface to a late subscriber, and a publish
  survives the root churning out.
\item
  \textbf{Metrics rebuilt (v4.10.1).}
  \texttt{rootedTopics()}/\texttt{peer.metrics()} were silently dead
  across early 4.x; rebuilt with two new fields ---
  \textbf{\texttt{seq}} (dense message counter) and
  \textbf{\texttt{cohortSize}} --- and cohort-aware aggregation (see
  §4.3).
\item
  \textbf{Cold-publish burst (v4.11.0).} A freshly-joined node's first
  publishes are automatically re-sent a few times over the first second,
  so a cold-start publish isn't lost while the routing table warms. No
  app action --- \texttt{pub} is unchanged.
\item
  \textbf{Nearest-replica reads (v4.11.1+).} \texttt{pull} is answered
  by the first replica the request reaches instead of always the root:
  exact for \texttt{pull(msgId)}, and \emph{recent}
  (eventually-consistent) for pull-latest, which spreads a hot read path
  off the root (see §4.3).
\end{itemize}

\subsubsection{What changed since v2.x}\label{what-changed-since-v2.x}

The pub/sub and identity surfaces were redesigned in the v3 line. If you
are migrating, the load-bearing differences are:

\begin{itemize}
\tightlist
\item
  \textbf{Two identity factories.}
  \texttt{createNodeIdentity(\{\ lat,\ lng\ \})} mints the
  \emph{connection} key (the nodeId); \texttt{createAuthorIdentity()}
  mints a location-free \emph{authorship} key (the Author ID). The old
  single \texttt{deriveIdentity()} is gone.
\item
  \textbf{Topics are descriptors, not strings.} A topic is
  \texttt{\{\ region?,\ owner?,\ name,\ write?\ \}}. \texttt{peer.pub}
  takes the descriptor; \texttt{peer.sub} / \texttt{pull} /
  \texttt{metrics} take the descriptor \textbf{or} a 66-hex topic ID.
  There is no \texttt{publisher}/\texttt{publishId} argument anywhere.
\item
  \textbf{Signing is explicit per publish.}
  \texttt{peer.pub(topic,\ message,\ \{\ signWith\ \})} requires a
  signer --- an author identity or the \texttt{ANONYMOUS} sentinel.
  There is no default author, the node key never signs publishes, and
  there is no \texttt{sign:\ false}.
\item
  \textbf{The envelope carries the topic descriptor.} A signed publish
  binds the exact \texttt{\{\ region,\ owner,\ name,\ write\ \}}, so a
  storing node recomputes the topic ID and enforces the write policy
  (\texttt{signer\ ===\ owner} for owned topics).
\end{itemize}

\begin{center}\rule{0.5\linewidth}{0.5pt}\end{center}

\subsection{A word on stability}\label{a-word-on-stability}

\texttt{@axona/protocol} 4.x is stable for the application surface
(§§1--9). The transport + protocol surface (§§10--16) is stable but
expect occasional additions. Low-level utilities (§§17--23) may grow but
won't change incompatibly.

For changelog and migration notes, see
\url{https://github.com/axona-net/axona-protocol/releases}.

\end{document}
